\documentclass[11pt,letterpaper]{article}
\usepackage[T1]{fontenc}
\usepackage{lmodern,microtype,amsmath,amssymb,amsthm,booktabs,array,longtable}
\usepackage[margin=0.92in]{geometry}
\usepackage[dvipsnames]{xcolor}
\usepackage{hyperref}
\hypersetup{colorlinks=true,linkcolor=MidnightBlue,citecolor=MidnightBlue,urlcolor=MidnightBlue,bookmarksnumbered=true,pdfauthor={Chatgpt},pdftitle={Planar Newtonian central configurations with five equal masses: finiteness for generic mass ratio}}
\urlstyle{same}
\setcounter{tocdepth}{1}
\setcounter{secnumdepth}{2}
\newtheorem{theorem}{Theorem}[section]
\newtheorem{lemma}[theorem]{Lemma}
\newtheorem{proposition}[theorem]{Proposition}
\newtheorem{corollary}[theorem]{Corollary}
\theoremstyle{definition}
\newtheorem{definition}[theorem]{Definition}
\newtheorem{remark}[theorem]{Remark}
\newcommand{\Q}{\mathbb{Q}}
\newcommand{\QQ}{\mathbb{Q}}
\newcommand{\R}{\mathbb{R}}
\newcommand{\RR}{\mathbb{R}}
\newcommand{\C}{\mathbb{C}}
\newcommand{\CC}{\mathbb{C}}
\newcommand{\F}{\mathbb{F}}
\newcommand{\Z}{\mathbb{Z}}
\newcommand{\KC}{\mathrm{KC}}
\newcommand{\suppfile}{\href{https://github.com/rongyulin/physics-open-problems-AI/raw/refs/heads/main/problems/planar-central-configurations/papers/FiveEqualFiniteness-checks.zip}{verification archive}}
\title{Planar Newtonian central configurations with five equal masses:\\ finiteness for generic mass ratio}
\author{Chatgpt}
\date{27 November 2026}
\begin{document}
\maketitle
\begin{abstract}
We consider the planar Newtonian six-body problem with masses proportional to $(1,1,1,1,1,t)$. We give a computer-assisted proof that, outside a finite set of positive algebraic ratios $t$, there are only finitely many real central configurations up to similarity. The five-equal-mass line lies entirely in the exceptional set of the known general-mass criterion. We specialize the catalogue of singular places on this line, retaining rather than discarding the leading systems whose ambient mass obstruction becomes identically zero. A cubic-order torque test treats the exceptional pair of triangles. Local force equations and real energy bounds then reduce the proof to two types of end. For one type the complete four-body force equations and the cross torque are incompatible on a projective special fibre in characteristic 65521. For the other, we classify or lift the leading solutions according to the position of the marked mass. We give the project-specific arguments and arithmetic tests in one text. The theorem allows finitely many undetermined mass ratios and does not assert finiteness of every signed complex fibre.
\end{abstract}

\section{Statement and structure of the proof}
Let $q_i\in\R^2$ denote the position of body $i$ and let $c$ be the center of mass. A collision-free configuration is \emph{central} if a common $\lambda$ satisfies
\begin{equation}\label{intro:phys}
 \sum_{j\ne i}m_j\frac{q_i-q_j}{|q_i-q_j|^3}=\lambda(q_i-c)
 \qquad (1\le i\le6).
\end{equation}
Such shapes include those rotating at a common angular speed as a relative equilibrium. A similarity here consists of a translation, a rotation and a positive rescaling. Finiteness with labelled bodies and this convention implies finiteness after identifying reflections or permutations of equal masses.
\begin{theorem}[Five equal masses]\label{intro:main}
There exists a finite set $S\subset\overline{\Q}\cap\R_{>0}$ such that, for every fixed $t\in\R_{>0}\setminus S$, equation~\eqref{intro:phys} with masses $(1,1,1,1,1,t)$ has finitely many real planar shapes up to similarity.
\end{theorem}
Multiplying all six masses by one positive number changes the force multiplier and preserves the shapes. Thus $t$ denotes a ratio. An infinite family of distinct real shapes at a single fixed ratio would force that ratio into $S$. The same inference cannot be drawn from one shape at each of infinitely many different mass ratios.

Smale's finiteness question asks for this conclusion at every fixed positive mass vector \cite{Smale}. For six bodies Moczurad and Zgliczy\'nski enumerated the planar configurations at the equal mass point and obtained local counts for masses sufficiently close to equality \cite{MZplanar,MZnear}. Dias and Pan proved generic finiteness for cross configurations with four bodies on an axis of reflection \cite{DP}. Tomar's general-mass theorem gives finiteness outside a proper algebraic subset of the six-mass space; its stated exceptional subset contains every vector with at least three equal entries \cite[Theorem 1.1 and Remark 6.14]{Tomar}. Chen and Tong give necessary mass conditions for a curve through a positive physical point \cite[Thm.~1.1, Cor.~1.2]{CT}. Aquino-Michaels states an ambient generic finiteness theorem for every number of bodies \cite[Thms.~1.1--1.3]{GenericN}. Its explicit mass specialization requires strictly separated valuations, whereas the five equal masses have identical valuations. A polynomial may be nonzero in six independent mass variables yet vanish identically on this entire line. Our task is to exhibit obstructions that survive the substitution of the five equal masses. The result does not compute $S$, nor does it require any positive parameter to be exceptional.

\paragraph{Proof strategy.}
Starting with infinitely many real shapes at a fixed mass, we produce a curve of complexified configurations. At a suitable place of the curve some coordinate or force diverges after balancing its two complex projections. Only a finite catalogue of patterns for these orders is possible. The issues are (i) which algebraic exclusions survive on the mass line, (ii) which of the remaining ends could belong to a curve obtained from physical configurations, and (iii) whether a solution of leading equations actually satisfies the full six-body equations. Each question requires a different test.
\begin{center}\small
\begin{tabular}{@{}>{\raggedright\arraybackslash}p{.21\linewidth}p{.71\linewidth}@{}}
\toprule
Input or case & What must be proved on the mass line\\\midrule
Diagram catalogue & Read the mass factors in the published identities with all six possible choices of the body of mass $t$. Deal explicitly with cells in which those identities supply no obstruction.\\[2pt]
Two triangles (98) & A leading-order cancellation can persist on the line. Compute higher orders and exclude a branch of the full signed equations.\\[2pt]
Diagram 93 and a marked binary in 88 & Carry the fixed-mass value of the potential and its marked-mass derivative from real points to a complex end. Check the bounds on the end values.\\[2pt]
Diagram 94 & The shared body's force in one projection is missing from the internal system. Resolve the three choices of marked body and the possible lifts.\\[2pt]
A marked quartet in 88 & Specialize the entire four-body force and torque system, including collapsed or singular shapes, and test the resulting algebraic fibres.\\
\bottomrule
\end{tabular}
\end{center}
We use Tomar's catalogue and local equations, with the diagram numbering of Chang and Chen. We derive the restrictions specific to the mass line below. Large catalogue identities, univariate data and executable arithmetic replays are distributed separately from the exposition. The \suppfile{} identifies inputs, commands and the coverage of each calculation. Arguments over algebraic closures are indicated explicitly wherever a calculation uses a finite field.

\begingroup
\small
\makeatletter
\renewcommand*{\l@section}{\@dottedtocline{1}{0em}{1.8em}}
\makeatother
\tableofcontents
\endgroup

\section{From real configurations to places of algebraic curves}\label{intro:places}
\subsection{Normalization, signs and the shape slice}
Translate so that $c=0$. Write $U=\sum_{i<j}m_i m_j/|q_i-q_j|$ and $I=\sum_i m_i |q_i|^2$. Contracting \eqref{intro:phys} with $m_iq_i$ and summing gives $\lambda I=U>0$. Rescaling all positions by $\lambda^{1/3}$ sets the multiplier to one. For real positions put $z_i=x_i+\mathrm i y_i$ and $w_i=x_i-\mathrm i y_i$. Write $z_{ij}=z_i-z_j$, $w_{ij}=w_i-w_j$ and $p_{ij}=p_{ji}$. We complexify the normalized equations by treating $z$ and $w$ as independent coordinates and imposing
\begin{equation}\label{intro:signed}
 z_i=\sum_{j\ne i}m_jp_{ij}^3z_{ij},\quad
 w_i=\sum_{j\ne i}m_jp_{ij}^3w_{ij},\quad
 p_{ij}^2 z_{ij}w_{ij}=1 .
\end{equation}
On the physical locus $p_{ij}=1/|q_i-q_j|$ is positive. Both roots of the last equation have to be allowed in the algebraic arguments. All pair differences and all inverse distances are units in the coordinate ring. Summing the first two equations with weights $m_i$ proves the two center-of-mass equations. Set $M=\sum m_i$. In this notation the extensions of $U,I$ satisfy
\begin{equation}\label{intro:energy}
 U=\sum_{i<j}m_im_jp_{ij},\qquad
 I=\sum_i m_i z_i w_i=\frac{1}{M}\sum_{i<j}m_i m_j z_{ij}w_{ij},\qquad I=U .
\end{equation}
For the last equality, add the contractions of the two force equations with $m_i w_i/2$ and $m_i z_i/2$. Each unordered pair contributes $m_i m_jp_{ij}^3z_{ij}w_{ij}=m_i m_jp_{ij}$.

The algebraic action $(z,w)\mapsto(bz,b^{-1}w)$ preserves all the equations and distances. Fix the slice $z_1-z_6=1$ to remove it, possible since this difference is a unit. The 15 squared distances have finite fibres on the slice. To see this put $u_i=z_i-z_6$, $v_i=w_i-w_6$, $D_{ij}=z_{ij}w_{ij}$. Then $v_1=D_{16}\ne0$, $u_i v_i=D_{i6}$, and
\[
 D_{16}u_i+v_i=D_{i6}+D_{16}-D_{i1}\qquad (2\le i\le5).
\]
Each $u_i$ satisfies a quadratic with nonzero leading coefficient. These choices determine the $v_i$, centering recovers $z_6,w_6$, and each $p_{ij}$ has only two choices. On the real locus the mutual distances determine the shape up to isometry, with at most two choices if only orientation-preserving isometries are identified. Consequently infinitely many normalized physical shapes give an infinite set of points of the slice with varying squared distances.

\subsection{Why a curve supplies a finite search}
For the constancy of the normalized energy at fixed masses see Albouy and Kaloshin \cite[Lemma~2]{AK}. Chen and Tong \cite[Lemmas~2.1--2.2]{CT} pass from a positive real arc to a complex curve of normalized configurations at that energy, with a varying squared distance. The Zariski closure of an infinite set in this affine variety has positive dimension. Choose an irreducible component on which some $D_{ij}$ is nonconstant. Regular functions constant on this component retain their values on any curve cut out in it. More explicitly, at a smooth point where the differential of a nonconstant distance is nonzero, general hyperplanes through the point cut out a curve on which the distance still varies. When physical inequalities or a specified body at the origin are required, first partition the real set and then take the relevant closure. The argument below derives the constants in \eqref{intro:energy} and a second scalar on generic fixed-mass fibres, so that their values at the chosen real stratum carry over to the curve.

Let $C$ be such a curve and $\bar C$ its smooth projective model. At a place $P$ write $v$ for the order of vanishing of a rational function, so poles have negative order. The pair-force terms in the two projections are $F^z_{ij}=p_{ij}^3 z_{ij}$ and $F^w_{ij}=p_{ij}^3 w_{ij}$. To compare their orders one may replace $(z,w)$ by $(bz,b^{-1}w)$ over a finite extension of the local parameter field. This makes the minimum order among the positions and pair forces in each projection the same. If that common minimum is negative we have a singular place. At a nonsingular place all the $D_{ij}$ have order zero: after balancing, each difference has nonnegative order and the two forces, with doubled orders $-v(z_{ij})-3v(w_{ij})$ and $-3v(z_{ij})-v(w_{ij})$, have nonnegative order as well. Hence both difference orders are zero.

There is a singular place, since at least one $D_{ij}$ has a pole on $\bar C$. For the nonzero rational function $\Delta=\prod_{i<j}D_{ij}$ the sum of its orders on $\bar C$ is zero. The contributions at nonsingular places vanish. Thus at least one singular place has $v(\Delta)\le0$, including the possibility that every contribution vanishes. This is the nonpositive-place criterion of \cite[Corollary 2.7]{Tomar}. It is a statement about a curve at \emph{fixed} masses. In particular it does not assume that those masses avoid the ambient exceptional set.

Scale the balanced orders so that the common minimum equals $-2$. A circle of a given colour marks a body whose coordinate has that order; a stroke of that colour joins two bodies whose pair-force term does. Orderings of the remaining differences and the sets of monomials attaining the least order in each equation refine a diagram into cells. Terms of a vanishing sum cannot have a unique lowest order unless their coefficient at that order vanishes. For example, write $a=v(z_{ij})$, $b=v(w_{ij})$, in any common scale. The order of $p_{ij}$ is $-(a+b)/2$ and that of $F^z_{ij}$ is $-(a+3b)/2$. These relations give linear constraints on orders and, at a tie, polynomial equations on leading coefficients. The catalogue makes the order bookkeeping finite. A surviving cell only imposes necessary equations on an end.

\begin{figure}[ht]
\centering
\begingroup
\setlength{\unitlength}{0.91pt}
\def\zstroke#1#2#3{\textcolor{MidnightBlue}{\qbezier#1#2#3}}
\def\wstroke#1#2#3{\textcolor{Mahogany}{\linethickness{0.75pt}\qbezier[20]#1#2#3}}
\def\vertex#1#2#3{\put#1{\textcolor{white}{\circle*{11}}}\put#1{#3}\put#1{\makebox(0,0){\scriptsize #2}}}
\def\zcircle{\textcolor{MidnightBlue}{\circle{12}}}
\def\wcircle{\textcolor{Mahogany}{\circle{12}}}
\begin{tabular}{@{}ccc@{}}
\begin{picture}(118,66)%
\zstroke{(9,10)}{(30.5,10)}{(52,10)}%
\zstroke{(9,10)}{(19.5,34.5)}{(30,59)}%
\zstroke{(52,10)}{(41,34.5)}{(30,59)}%
\zstroke{(9,10)}{(19.5,18.5)}{(30,27)}%
\zstroke{(52,10)}{(41,18.5)}{(30,27)}%
\zstroke{(30,59)}{(30,43)}{(30,27)}%
\wstroke{(82,43)}{(92.5,32)}{(103,21)}%
\vertex{(82,43)}{0}{\wcircle}%
\vertex{(9,10)}{1}{\zcircle}%
\vertex{(52,10)}{2}{\zcircle}%
\vertex{(30,59)}{3}{\zcircle}%
\vertex{(30,27)}{4}{\zcircle}%
\vertex{(103,21)}{5}{\wcircle}%
\end{picture} &
\begin{picture}(118,66)%
\zstroke{(9,10)}{(31.5,10)}{(54,10)}%
\zstroke{(9,10)}{(20,32.5)}{(31,55)}%
\zstroke{(54,10)}{(42.5,32.5)}{(31,55)}%
\wstroke{(77,10)}{(92.5,10)}{(108,10)}%
\wstroke{(77,10)}{(84.5,31.5)}{(92,53)}%
\wstroke{(108,10)}{(100,31.5)}{(92,53)}%
\vertex{(77,10)}{0}{\wcircle}%
\vertex{(9,10)}{1}{}%
\vertex{(54,10)}{2}{\zcircle}%
\vertex{(31,55)}{3}{\zcircle}%
\vertex{(108,10)}{4}{\wcircle}%
\vertex{(92,53)}{5}{\wcircle}%
\end{picture} &
\begin{picture}(118,66)%
\zstroke{(9,10)}{(10,33.5)}{(11,57)}%
\zstroke{(9,10)}{(18.5,20)}{(28,30)}%
\zstroke{(11,57)}{(19.5,43.5)}{(28,30)}%
\wstroke{(107,50)}{(107,31)}{(107,12)}%
\wstroke{(107,50)}{(84.5,40.5)}{(62,31)}%
\zstroke{(9,10)}{(35.5,20.5)}{(62,31)}%
\zstroke{(11,57)}{(36.5,44)}{(62,31)}%
\zstroke{(28,30)}{(45,30.5)}{(62,31)}%
\wstroke{(107,12)}{(84.5,21.5)}{(62,31)}%
\vertex{(107,50)}{0}{\wcircle}%
\vertex{(9,10)}{1}{\zcircle}%
\vertex{(11,57)}{2}{\zcircle}%
\vertex{(28,30)}{3}{\zcircle}%
\vertex{(107,12)}{4}{\wcircle}%
\vertex{(62,31)}{5}{}%
\end{picture} \\
\small Diagram 88 & \small Diagram 93 & \small Diagram 94
\end{tabular}
\endgroup
\caption{Representative labelled diagrams used in the last reductions. Solid and dotted edges denote leading forces in the $z$ and $w$ projections, respectively. Rings in the corresponding colours mark leading coordinates. The shared vertex in diagram 94 is $5$ in this catalogue labelling. Each later force calculation relabels its vertices locally. The diagrams record orders at a place; the ordinary six-body equations still contain every pair interaction.}\label{intro:graphs}
\end{figure}

\subsection{Using a generic parameter and returning to a fixed mass}\label{intro:generic}
We use two elementary algebraic facts. First, equations and inequations with rational coefficients define a constructible set. An inequation $d\ne0$ can equivalently be represented as $hd=1$ using one extra variable $h$. The projection of a constructible set to the affine $t$-line is constructible. An infinite constructible subset of an algebraic line over $\overline{\Q}$ contains a nonempty open set. Thus emptiness over $\overline{\Q(t)}$ implies that the possible parameter values lie in a finite algebraic set; the same conclusion holds after base change to $\C$. Radical choices are retained as polynomial variables subject to their equations.

Second, testing the equations at a convenient single mass rarely proves emptiness over $\overline{\Q(t)}$ by itself. The curve $(t-2)x=1$, for example, has a point in the generic fibre and none in the affine fibre $t=2$. Two ways to prevent such a loss occur below. In the first, a finite map has no missing branches over the chosen regular point, and a unit denominator and a nonzero residue witness failure on every sheet. In the second, one takes a valuation of an alleged generic solution, controls its projective limit with the distance and force equations, and tests also the boundary shapes for which an affine normalization breaks down. The second issue is essential for diagram 88, where a shrinking cluster can have zero multiplier in its rescaled limit. Real inequalities enter later only on real configurations or on scalar quantities inherited from them.
\section{Specializing the arbitrary-mass catalogue}\label{early:census}

We work with the finite list of nonpositive cells in the catalogue of
\cite[Props.~3.1, 4.5 and 5.3]{Tomar}. Its two coordinate trees record orders
of differences, including differences with the distinguished origin. At a node,
the projections of distinct children have distinct leading coefficients. For
each body equation the cell records all terms attaining its minimum valuation.
These discrete data suffice for the tests below. We use the diagram numbers of
Chang and Chen~\cite{CC}. A placement is a triple $(c,i,s)$, where $c$ is
the catalogue class, $i$ the zero-based cell number within that class, and
$s\in\{1,\ldots,6\}$ the position of the mass $t$. This convention keeps the
six distinguished slots even in cells with automorphisms. All counts below use
the data for Tomar's Proposition~5.3, including its second search for cores. Here a core is a subsystem of the initial equations used to certify a necessary polynomial identity in the masses.

\paragraph{Specialization with a fixed mass.}
Of 91,273 cells, 52,432 have a single term of minimum valuation in one
necessary identity. Such a term has nonzero leading coefficient at positive
masses, by \cite[Lemma~5.1]{Tomar}. For 37,808 further cells there is an
identity $\sum_i c_i f_i=h(m)$ in the initial equations; 990 more have such
an identity after fixing a mass to one. Evaluating at a place annihilates
the left-hand side. These 38,798 cells use 1,928 distinct certificates;
43 cells have no assigned identity. The distinction is material because
a nonzero polynomial in six masses can vanish identically on this line.

The specialization rule is as follows. For a term $a_\alpha m^\alpha$ of $h$,
accumulate $a_\alpha$ in a dictionary with integer key
\[
e(\alpha,s)=
  \begin{cases}
   \alpha_s-\sum_j\alpha_j & \text{if the certificate fixed slot }s,\\
   \alpha_s & \text{otherwise}.
  \end{cases}
\]
Indeed, fixing slot $f$ divides every physical mass by its original
$m_f$ before applying the identity. The accumulated coefficients are those
of a Laurent polynomial in $t$. If any sum is nonzero, this placement is
impossible off the finite root set of that polynomial. In the 232,788
tests, 203,119 give a nonzero polynomial and 29,669 give zero. In 30,475
cells all six slots give nonzero polynomials. For example, class 29 cell
1617 (diagram 8) has certificate
$m_6(m_2^6-m_5^6)(m_3^6-m_4^6)$, which vanishes for all six placements;
in diagram 63 the certificate $m_1+m_2$ excludes all positive masses.

\paragraph{Triangles and four-body cliques.}
Call three labels a triangle if in each tree they occupy three different
children of one node. A vertex is balanced in a coordinate when the two
forces from the other triangle vertices are exactly the minimizing terms in
its equation there. When at least two vertices are balanced in one coordinate
and at least one in the other, \cite[Cor.~9.10]{Tomar} gives
\[
Q(a,b,c)=(ab)^2+(bc)^2+(ca)^2-2abc(a+b+c)=0 .
\]
Its values on $(1,1,1)$ and $(1,1,t)$ are $-3$ and $1-4t$.
Testing every triple in both coordinate orders finds 86,875 qualifying cells
in the full catalogue. In the 29,669 placements presently under consideration,
23,349 contain a qualifying triangle avoiding $s$ and another 4,788 have
qualifying triangles containing $s$. Thus 1,532 remain generically. Of these,
640 meet one of the positive-mass exclusions in \cite[Lemmas~9.2--9.7]{Tomar}.
For this last finite join one may use the published cell array
\texttt{tiers\_v2.jsonl}: select a row when its
\texttt{cell\_level\_allmass} intersects
$\{\mathtt{T,P,C,B0,K,L}\}$ and identify the cell by $(c,i)$.
This leaves 892 placements.

For completeness we specify the next incidence test. Enumerate triples and
quadruples $K$ whose edges are all double, or all of one colour with no
circle of that colour on $K$. In each coordinate require every internal
force of $K$ to be minimal at its target. Let $a$ count vertices of $K$
with an external minimizing force and let $b$ be one if any position is
minimal, zero otherwise. The forcing dimension is $d=\min(|K|,a+b)$:
the external forces span individual coordinate vectors, and the position
adds at most the all-ones vector. For a triple require
$\max(d_z,d_w)\le2$ and $d_z+d_w\le3$; for a quadruple require
$\max(d_z,d_w)\le1$. Corollary~9.18 of~\cite{Tomar} applies. In the
triple case its two possible mass relations from Theorem~9.9 give $Q=0$
or $m_s^2(m_a+m_b)^4=m_a^3m_b^3$ for some ordering here denoted $(a,b,s)$.
Both fail at three unit masses. In the quadruple case the closed exceptional
set $B_E$ misses the unit quadruple by Proposition~7.8. Each pullback
therefore has finitely many points. This test selects 86,938 cells of the
full input and removes 28 further placements, all in diagram 31.

In diagram 8 one can sharpen the generic triangle argument. Its 1,866 cells
include 1,842 with a balanced triangle. Enumerate all triples and both
coordinates on the remaining 24. In each there are two balanced vertices
$a,b$ in one coordinate; in the other, the minima for each are precisely
its position and its two triangle forces, and the least common ancestor of
$a,b$ lies strictly below their common meeting with the origin. These are
direct tests of set equality and ancestry. The internal forces on the
third vertex sum to zero in the first coordinate by mass-weighted
cancellation. The leading positions of $a,b$ in the other coordinate
agree and are nonzero by the ancestor test. The equal-force alternative of
\cite[Thm.~9.9(ii)]{Tomar} gives the second mass relation just displayed.
For $m_s=t$ it gives $16t^2=1$; for $m_a=t$ or $m_b=t$ it gives
$(1+t)^4=t^3$, impossible since $(1+t)^4/t^3=t+4+6/t+4/t^2+1/t^3>1$.
All unit masses also fail. Thus only $t=1/4$ can supply a nonpositive
end in diagram 8. This sharper assertion is unnecessary for the later count.

\paragraph{Whole-diagram exclusions.}
Earlier reductions by Aquino-Michaels give mass obstructions for diagrams 52, 63, 65 and 67 and use positive inertia to exclude diagram 6 on branches arising from real configurations \cite{Residual}. Here each exclusion used from the catalogue is checked on the specialized mass line.
Propositions~6.4 and 6.7 and Theorem~6.5 of~\cite{Tomar} rule out
diagrams 90, 82 and 89 respectively. Proposition~6.8 requires three
disjoint equal-mass pairs for 81, impossible for $t\ne1$.
For 64 and 66, Theorems~7.12--7.13 require a specified quadruple to
belong respectively to $B_E$ and $B_C$. These closed algebraic sets
both miss $(1,1,1,1)$, by Propositions~7.8 and 7.11 and the discussion
following Theorem~7.13. Consequently their pullbacks to each labelled
mass line are proper closed subsets. The next section establishes the
additional whole-diagram exclusion needed for 98.

\section{The two-triangle degeneration (diagram 98)}\label{early:98}

Consider masses proportional to \( (1,1,1,1,1,k) \), with \(k>0\). In a diagram-98 singular place, two triples separate in opposite complex coordinates. After interchanging the triples and relabelling bodies, their masses are \(\mu=(1,1,1)\) and \(\nu=(1,1,k)\). We first restrict their leading coefficients, then use a cubic-order necessary condition to exclude the remaining shapes for generic mass. In this section we write $k=t$ to leave a separate variable available for a triangle shape. For arbitrary positive masses Chen and Tong already exclude a first limit consisting of two Lagrange triples on a curve through a positive physical point \cite[Prop.~5.21]{CT}. For first limits containing an Euler triple they derive the nine-term signed cross-torque equation \cite[eq.~(113)]{CT}. We use the five-equal-mass hypothesis to analyze higher-order lifting for signed curves.

\begin{proposition}\label{early:leading} There exists an explicitly defined nonzero polynomial \(P\in\mathbb Z[k]\) with \(P(3)\ne0\) such that, for every \(k>0\) with \(P(k)\ne0\), every solution to the diagram-98 leading shape and cross-torque conditions has both of the following properties.

\begin{enumerate}
\item The projected shape \(x\) of the all-equal triple satisfies \(x^2-x+1=0\).
\item The cross sign matrix is of rank one over the complex numbers. Its signs can be written as \(\epsilon_{jl}=\delta_j\eta_l\), independently of the chosen square-root branches.
\end{enumerate}

\end{proposition}

The conclusion holds at \(k=3\). It covers every labelled allocation of the five interchangeable unit masses. The projected complex shape need not represent a real equilateral triangle.

\subsection*{Input from the central-configuration equations}

By \cite[Lemmas~6.9 and 6.11, Thm.~6.12]{Tomar} the following conditions are necessary. For the triple with masses \(m=(1,1,h)\) and shape \(z\), set

\[
(\ell_1,\ell_2,\ell_3)=(1+hz,hz-1-h,1-2z),\qquad
\Lambda_j=\prod_{i\ne j}\ell_i.
\]

All three linear forms are nonzero at a diagram-98 place. Write \(S_h\) for Tomar's degree-34 shape polynomial. With \(S_1(x)=S_k(y)=0\), the cross-torque vanishes only if some signed sum

\[
\sum_{j,l=1}^{3}\epsilon_{jl} r_j s_l = 0,\qquad
r_j^2=\mu_j^2\Lambda_j(x;1),\quad
s_l^2=\nu_l^2\Lambda_l(y;k)
\]

vanishes. The signs of a row or column change when a corresponding square-root branch changes. Rank and the equation with the transformed sign matrix are invariant. Overall sign is irrelevant, so require \(\epsilon_{11}=1\). Tomar's Remark 6.14 observes that rank-one matrices include the factorable sums that vanish for an equal-mass triple. The corresponding original obstruction polynomial therefore vanishes identically on this line. The proposition tests the other signs and the other equal-triple shapes separately.

\subsection*{Reduction of the shape equations}

Put \(E(z)=z^2-z+1\) and

\[
L=E(z)(z^4-5z^3+9z^2-5z+1)(z^4+z^3-2z+1)(z^4-2z^3+z+1).
\]

For the equal triple retain the degree-26 polynomial

\[
\begin{aligned}A=L{}&(2z^4-z^3+2z^2-2z+1)(z^4-2z^3+2z^2-z+2)\\
&(2z^4-7z^3+11z^2-7z+2).\end{aligned}
\]

For the other triple put

\[
\begin{aligned}
B_k=L{}&[2z^5-5z^4+4z^3-(3k+3)z^2+(3k+2)z-(k+1)]\\
&[z^4-2z^3+(k+1)z^2-kz+k+1]\\
&[2z^5-5z^4+4z^3+(3k+1)z^2-(3k+2)z+k+1]\,
[z^2-z-k-1].
\end{aligned}
\]

Its degree is 30 for every \(k\). Direct multiplication in \(\mathbb Z[k,z]\) gives

\[
\begin{aligned}S_1&=A(z)(2z-1)^2 E(z)(z+1)^2(z-2)^2,\\
S_k&=B_k(z)(2z-1)^2 E(z).\end{aligned}
\]

In both products the retained polynomial already contains \(E\). Every other removed root makes at least one \(\ell_i\) vanish. Thus these equations retain every admissible shape even at special masses. The equal-triple polynomial \(A/E\) has no common root with \(E\); in particular, a non-equilateral equal-triple branch cannot specialize to an equilateral root at 3. Both retained polynomials have constant leading coefficient 4; no shape escapes to infinity by a degree drop. To verify the product identities directly, take
\[
 S_h=L\prod_{a,b\in\{1,-1\}}
 [2(z-1)^3z^2+(z-1)^2z^2+ah(z-1)^3-bhz^3-a(z-1)^2-bz^2].
\]
This is the three-body shape polynomial in Tomar's chart; ordinary polynomial multiplication gives the displayed factorizations.

\subsection*{An integer obstruction polynomial}

Let \(\mathcal N\) be the 240 canonical matrices of rank greater than one and \(\mathcal F\) the 16 canonical matrices of rank one. For either class \(\mathcal C\), define

\[
H_{\mathcal C}(r,s)=
\prod_{e\in\mathcal C}\left(\sum_{j,l} e_{jl}r_js_l\right).
\]

Changing the sign of any \(r_j\) or \(s_l\) permutes these factors up to overall signs. If \(e_{11}\) is affected, \emph{every} factor needs a minus sign to restore the normalization, and the number of factors is even. Hence every exponent of every radical in the product is even. Substituting the six square identities above yields \(h_{\mathcal C}(x,y,k)\in\mathbb Z[x,y,k]\). Define

\[
P(k)=\operatorname{Res}_x\big(A,\operatorname{Res}_y(B_k,h_{\mathcal N})\big)
\operatorname{Res}_x\big(A/E,\operatorname{Res}_y(B_k,h_{\mathcal F})\big).
\]

All polynomial evaluations take place after replacing radicals by their squares in the invariant product. This is an explicit polynomial specified by products and determinants rather than by expanded coefficients. A root pair with either a nonfactorable sign matrix, or a non-equilateral equal-triple shape and a factorable matrix, makes the corresponding nested resultant zero.

It suffices to check \(P(3)\ne0\). At 3 both shape polynomials have their full degrees and all roots are simple. The product formula for a resultant of a fixed-degree outer polynomial multiplies the inner values over all outer roots, up to a nonzero power of the leading coefficient, even if the degree of the inner polynomial drops in a specialization. Thus the finite check below establishes \(P(3)\ne0\) without assuming that a component of any system of rational equations meets the tested affine chart. No inversion of a parameter, radical, or shape coordinate is needed to construct \(P\). For a sign class of size \(M\), \(H_{\mathcal C}\) is homogeneous of degree \(M/2\) in the squared coordinates of each triple separately. The substituted entries have shape degree at most two, and the mixed entries have parameter degree at most four. Hence \(h_{\mathcal C}\) has degrees at most \((M,M,2M)\) in \((x,y,k)\). The mixed-shape polynomial has parameter degree at most four and shape degree 30. The Sylvester determinant bounds the inner resultant's parameter degree by \(30(2M)+4M=64M\); the outer shape polynomials are constant in the parameter. In particular,

\[
\deg P \le 26\cdot64\cdot240+24\cdot64\cdot16=423936.
\]

The finite calculation below proves nonvanishing, and hence the finite-exception assertion.

\subsection*{The finite witness}

The computation at 3 uses \(a_j=\mu_j/\sqrt{\ell_j(x;1)}\) and \(b_l=\nu_l/\sqrt{\ell_l(y;3)}\). All linear forms are nonzero on both root sets, so multiplying by one nonzero common scalar for each triangle, with harmless changes of row/column signs, identifies this test with the signed sums in \(r,s\).

Set $B=2^{48}$, $R=2^{-39}$ and $H=2^{-32}$. The following table specifies all shape centres $c=(X+\mathrm{i}Y)/B$. Include the conjugate when $Y>0$. In the table $L_i$, $A_i$ and $F_i$ enumerate, from left to right, the four factors of $L$, the three additional factors of $A$, and the four additional factors of the mixed shape polynomial at $k=3$. The first four factors serve both triples.

\begin{center}\small
\begin{tabular}{crr}
Factor & $X$ & $Y$\\\hline
$L_{1}$ & 140737488355328 & 243764480361061\\
$L_{2}$ & 106407554119526 & 52841401652252\\
 & 597279887657114 & 296605882013313\\
$L_{3}$ & -315804910946458 & 296605882013312\\
 & 175067422591130 & 52841401652251\\
$L_{4}$ & -133295707497992 & 125192134512859\\
 & 414770684208648 & 125192134512859\\
$A_{1}$ & -79046552355778 & 293796310620219\\
 & 149415296533442 & 107618113454196\\
$A_{2}$ & -67657734116147 & 251466915075327\\
 & 349132710826803 & 251466915075327\\
$A_{3}$ & 132059680177214 & 107618113454196\\
 & 360521529066434 & 293796310620219\\
$F_{1}$ & -126374236140503 & 402377118999033\\
 & 144102185109101 & 94676061820442\\
 & 668231543839445 & 0\\
$F_{2}$ & -9532650419105 & 348735311877243\\
 & 291007627129761 & 348735311877242\\
$F_{3}$ & -386756567128790 & 0\\
 & 137372791601555 & 94676061820442\\
 & 407849212851159 & 402377118999033\\
$F_{4}$ & -439538041617825 & 0\\
 & 721013018328481 & 0\\
\end{tabular}
\end{center}

At each such centre $c$, choose the square root of $m_j^2/\ell_j(c)$ having positive real part, or nonnegative imaginary part on the imaginary axis. Define $w_j$ by rounding its real and imaginary parts separately to the nearest multiple of $1/B$ (either choice at a tie). This prescription determines all auxiliary radical centres. One can use rational bisection for the ordinary positive square roots in
$\sqrt{x+\mathrm{i}y}=\sqrt{(|x+\mathrm{i}y|+x)/2}
 +\mathrm{i}\,\operatorname{sgn}_+(y)\sqrt{(|x+\mathrm{i}y|-x)/2}$,
where $\operatorname{sgn}_+(0)=1$.
 For a proposed root centre \(c\), write the exact Taylor coefficients of its assigned factor as \(f(c+w)=\sum d_i w^i\). The verifier proves \(|d_1|R>|d_0|+\sum_{i\ge2}|d_i|R^i\). Rouché's theorem gives exactly one zero in each disk. Disjointness and the degrees certify completeness.

For the proposed radical centre \(w\), let \(\ell_0=\ell(c)\) and let \(s\) be the slope of this linear form. The true value differs from \(\ell_0\) by at most \(|s|R\). The verifier proves

\[
2|\ell_0w|H>
|\ell_0w^2-m_j^2|+|\ell_0|H^2+|s|R(|w|+H)^2.
\]

A second application of Rouché's theorem isolates a radical root in each disk. The other root is its negative. The global sign enumeration covers all choices. Given centres \(w_j,v_l\) with error at most \(H\), the product \(a_jb_l\) has error at most \((|w_j|+|v_l|)H+H^2\). The verifier sums these errors and compares them to a lower bound on the modulus of every signed sum. It bounds \(|z|\) below by \(\max(|\Re z|,|\Im z|)\) and above by \(|\Re z|+|\Im z|\), so every comparison uses integers or exact rational arithmetic.

There are 26 equal-triple roots and 30 mixed-triple roots. All \(26\cdot30\cdot240=187200\) rank-greater-than-one combinations are excluded. All \(24\cdot30\cdot16=11520\) non-equilateral rank-one combinations are also excluded. The least certified lower bound on the modulus over these two classes exceeds 0.036477197 in the inverse-root coordinates used for this finite witness. Among 199680 total combinations, 240 fall within the interval error. Every one has an equilateral equal-triple root and factorable signs. This count is a count of signed algebraic tests, not of central configurations.

For either equilateral root, set \(Y_j=\mu_j^2\Lambda_j(x;1)\). The product of the four three-term signed sums equals

\[
Y_1^2+Y_2^2+Y_3^2-2(Y_1Y_2+Y_2Y_3+Y_3Y_1)=9E(x)^2.
\]

It follows that for both equilateral roots some signed sum over the equal triangle vanishes. Multiplying this identity by an arbitrary signed sum over the mixed triangle produces a factorable solution of the scalar torque equation. Additional information from the governing equations is required to decide whether any such leading solution extends to a singular place.

\subsection*{A force identity on the surviving branch}

There is a useful further restriction in the physical leading coordinates. Suppose
\(k>0\), \(P(k)\ne0\), and \(\Delta(k)\ne0\), where $\Delta$ is specified below. Write \(Z_j\) for the leading
projected coordinate of equal body \(j\), and \(W'_l\) for that of mixed
body \(l\), in the two opposite coordinates of Tomar's Lemma 6.11. Then
there are nonzero \(\gamma_l\) such that for the leading cross inverse
distances
\[
   \widehat p_{jl}=\gamma_l Z_j .
\]
In particular, for each mixed body \(l\), the three values
\(\widehat p_{jl}^3\) are equal, and
\[
\sum_{j=1}^3\widehat p_{jl}=0,\qquad
\sum_{j=1}^3\widehat p_{jl}^3 Z_j=0.
\]

To see this, first rule out a cancellation in the mixed triple.
At its shape \(y\), set \(Y_l=\nu_l^2\Lambda_l(y;k)\) and
\(D=\sum_lY_l^2-2\sum_{l<q}Y_lY_q\). Let $\Delta(k)=\operatorname{Res}_y(B_k,D)$. Nonvanishing at $k=3$ suffices here. In this specialization
\[
D=8505y^4-17010y^3+2395y^2+6110y+1105 ;
\]
reducing $D$ and the displayed $B_3$ modulo $101$ and running the Euclidean algorithm gives monic gcd $1$. Their leading coefficients are nonzero modulo $101$.
Thus the mixed factor of a rank-one
vanishing sum is nonzero, so the equal factor must vanish.

The three equal masses have centered coordinates proportional, in some
order, to \((1,\omega,\omega^2)\), where \(\omega\) is a primitive cube
root of unity. This follows by substituting \(x^2-x+1=0\) into the
linear forms above. The three numbers \(Z_j^3\) have a common nonzero
value, so for a common constant \(c\) one possible set of radicals
of \(1/Z_j\) is \(cZ_j\). These sum to zero. A signed zero sum of
three nonzero complex numbers has only one sign pattern up to a
global sign: two inequivalent patterns would isolate a single nonzero
term after adding or subtracting the relations. The equal-triple
radicals in a vanishing factor must therefore be proportional to
\(Z_j\). Factorability and
\(\widehat p_{jl}^{\,2}Z_jW'_l=-1\) give the displayed formula.
Both moment identities follow from \(\sum_j Z_j=0\).

\subsection{The cubic-order exclusion}\label{early:98cubic}

\begin{proposition}\label{early:cubic-prop} Outside a finite algebraic set of parameters
\(k>0\), no diagram-98 singular place exists for the masses
\((1,1,1,1,1,k)\). The exclusion holds at \(k=3\).\end{proposition}

\paragraph{A chart near the two triangles.}
Use Tomar's exponents in Lemma 6.10. By a complex rotation fix
\(z^A_{12}=1\), and set \(h^2=1/w^B_{12}\) in these rotated
coordinates. A further quadratic extension to extract \(h\) is harmless.
Write the physical positions in the chart as
\[
 (Z_j,W_j)\quad(j\in A),\qquad (h^2 Z'_l,h^{-2} W'_l)\quad(l\in B).
\]
Thus \(Z_{12}=W'_{12}=1\), \(h\to0\), and all the coordinate
vectors in capital letters have finite limits. The designated pair
differences and all the individual \(Z_j,W'_l\) at the limit are nonzero.
The two triples obey their own isolated normalized three-body equations at
the limit. The two triple centres are included as chart variables; stationarity of the harmonic translation terms recovers the centre-of-mass condition. For \(h\ne0\),
the chart gives configurations modulo the common rotation, locally up
to the choice of sign of \(h\).

With each local inverse-distance branch chosen by the given place, the
normalized central equations imply stationarity of
\(V=\frac12\sum_i m_i z_iw_i+\sum_{i<j}m_im_jp_{ij}\).
Write \(V_A,V_B\) for the analogous isolated-triple functions. Each cross
inverse distance, divided by \(h\), is locally analytic and equals
\[
 (-Z_jW'_l)^{-1/2}
 (1-h^2Z'_l/Z_j)^{-1/2}(1-h^2W_j/W'_l)^{-1/2},
\]
with its own leading sign. Thus, as an analytic function of all chart
coordinates,
\[
 V=V_A+V_B+h C_0+h^3 C_1+O(h^5).
\]
The preceding leading-order test gives a rank-one sign matrix, hence
\(C_0=a(Z)b(W')\), where the signed terms of \(a\) are branches of
\(1/\sqrt{Z_j}\) and those of \(b\) are branches of a common nonzero
constant times \(\nu_l/\sqrt{W'_l}\). At the limiting equal triple
\(a=0\), the signed radicals in \(a\) are \(c Z_j\), and
\(\partial a/\partial Z_j=-c/2\) for all three bodies. The mixed
triple has \(b\ne0\) when $\Delta(k)\ne0$, by the signed-cancellation
test above. Denote by \(g_{jl}\) the leading cross
radical at the limiting chart point. Then
\[
 C_1=\frac12\sum_{j,l}\nu_l g_{jl}
    \left(\frac{Z'_l}{Z_j}+\frac{W_j}{W'_l}\right)
 \quad\text{at that point}.
\]

\paragraph{Possible equal-triple lifts.}
Put \(x=Z_{13}\), \(v=W_{12}\), and \(x'=W_{13}/v\).
At the limit \(x=r\), where \(r^2-r+1=0\).
For a Lagrange-type solution of the three-body equations,
all three internal inverse distances have cube \(1/3\).
The ratios \(d=D_{13}/D_{12}\) and \(e=D_{23}/D_{12}\), with
\(D_{ij}=Z_{ij}W_{ij}\), satisfy \(d^3=e^3=1\),
\(r^2-(1+d-e)r+d=0\), and \(x'=d/r\).
Checking the three cube roots gives either \(d=e=1\),
\(x'=1/r\), or \(d=r^2,e=(r-1)^2\), \(x'=r\).
For an Euler-type lift, \(x'=r\). Here the distances are
\(v,vr^2,v(r-1)^2\), which all have the same cube. The three inverse
cubes therefore equal \(c\epsilon_{ij}\) for a common scalar \(c\ne0\)
and real signs \(\epsilon_{ij}\). The weighted graph Laplacian on the
three equal bodies with these edge signs is real symmetric, and the
centered vector \(Z\) is its eigenvector with eigenvalue \(1/c\).
This eigenvalue is real. Its conjugate vector is a second, independent
centered eigenvector, since the ratio \(Z_{13}/Z_{12}=r\) is nonreal.
The Laplacian is therefore scalar on the centered plane and zero on
the constant vector. Its three off-diagonal entries are equal, so all
three inverse cubes agree and the force equation gives their common
value \(1/3\).
Thus \(p_{ij}^3=1/3\) at both kinds of lift. Let
\(I_A=\sum_j Z_j W_j=(D_{12}+D_{13}+D_{23})/3\) at the centered
limit. If \(x'=1/r\), all three \(D_{ij}\) are equal and \(I_A^3=9\).
If \(x'=r\), \(I_A=0\) because the three \(D_{ij}\) sum to zero.

\paragraph{Regular equal-triple variables.}
In the independent centre variables \(z_c,w_c\), the translation part
of \(V_A\) is \(3z_cw_c/2\), with invertible Hessian. The centered
part is \(\sum_{ij} f_{ij}(D_{ij})\), where
\(f_{ij}(D)=D/6+D^{-1/2}\) with the chosen inverse-distance signs.
At both lifts \(f_{ij}'=0\) and \(f_{ij}''=1/(4D_{ij})\).
Indeed \(p^2D=1\) and \(p^3=1/3\).
The three distance variables are
\[
  (D_{12},D_{13},D_{23})=(v,vxx',v(x-1)(x'-1)).
\]
If \(x'\ne x\), their Jacobian in \((v,x,x')\) is invertible; hence
all five equal-triple coordinates in the rotation slice are regular.

At the collinear lift set \(p=(x+x')/2\), \(u=(x-x')/2\).
Here the last two distances are \(v(p^2-u^2)\) and
\(v((p-1)^2-u^2)\). In the \((v,p)\) Hessian block,
\(H_{vv}=\sum D_{ij}/(4v^2)=0\) and
\(H_{vp}=\sum(\partial_pD_{ij})/(4v)=(2r-1)/2\ne0\).
Thus the four variables \(z_c,w_c,v,p\) are regular even at this lift.
At \(h=0\) eliminating them gives analytic functions of \(u^2\), with
centres zero and \(p=r+O(u^2)\). All distance variations in this
solution are \(O(u^2)\). Since all \(f'_{ij}=0\), the reduced
equal-triple potential is constant up to \(O(u^4)\).

\paragraph{The critical value in the remaining variables.}
Treat the mixed-triple coordinates, denoted by \(Y\), as independent
parameters in a neighbourhood of their limiting values \(Y_0\).
Temporarily truncate the potential at
\(V^{\mathrm{tr}}=V_A+V_B+h a(Z)b(W')\). For every \(h,Y\) near
the limiting point, the limiting equal triangle with just its three
\(W_j\) translated by the same quantity
\(-2h b(W')\partial_{Z_1}a\) is exactly stationary in \emph{all}
equal-triple variables of the truncated problem. In fact the
translation leaves internal distances and all \(Z_j\) unchanged,
and the change of each \(Z_j\)-derivative of the harmonic term is half
the common translation. The value of \(V_A\) is unchanged since
\(\sum_j Z_j=0\). This observation holds in moving as well as fixed
mixed-triangle centres.

Use the implicit function theorem on the regular equal-triple variables
just identified, first for the truncation, then for \(V\). The full and
truncated critical variables differ by an analytic multiple of \(h^3\).
Thus the two reduced potentials differ by \(h^3 J\), where \(J\)
is analytic in \(h,Y\) and, in the collinear case, also \(u\).
At the limiting configuration \(J=C_1\). For example, this identity
follows by expanding the truncation about its stationary regular
variables, which gives only a quadratic error in their shift. This
argument can equally be applied in completed local power-series rings;
along a Puiseux branch the estimates below are estimates in its
uniformizing parameter. No Hessian of \(V_B\) is inverted.

In the noncollinear case the truncated critical value is
\(V_B(Y)+\mathrm{const}\) by the exact translation branch. The
partial derivative of the full critical value in \(h\), at fixed
\(Y\), consequently forces \(C_1=0\) in the limit.

In the collinear case the exact translation branch lies at \(u=0\);
it is critical also in \(u\) there. The reduced truncation therefore
has the analytic form
\[
  V_B(Y)+\mathrm{const}+u^2 L(h,u,Y),\qquad L(0,u,Y)=O(u^2).
\]
The coefficient \(d(Y)=\partial_h L(0,0,Y)\) is nonzero near
\(Y_0\). Indeed, differentiating the reduced truncation at \(h=0\)
gives precisely \(b(W')a(Z)\) at the regular critical variables there.
The centered \(Z\)-coordinates have shape \(x=p+u=r+u+O(u^2)\).
In these coordinates the vanishing equal-triple radical sum has an
\emph{exact double zero} at \(x=r\). The product of its four signed
choices is a nonvanishing analytic factor times \(9(x^2-x+1)^2\),
and the other three signed choices are nonzero by uniqueness of the
sign relation. This proves the claimed nonvanishing, since \(b(W')\ne0\).

Stationarity of the full reduced potential in \(h\) gives
\(0=u^2[d(Y)+O(h,u)]+O(h^2)\), hence \(u=O(h)\). Stationarity
in \(u\) then gives
\[
 0=2h d(Y)u+O(u^3+h u^2+h^2u+h^3),
\]
so \(u=O(h^2)\). Dividing the original \(h\)-derivative by \(h^2\)
now yields \(3J(0,0,Y_0)=0\), and again \(C_1=0\). These differentiations
are partial derivatives in the independent coordinates, using
stationarity, and do not require a rate estimate on \(Y-Y_0\).

\paragraph{The finite cubic torque test.}
The force identity above gives \(g_{jl}=c Z_j b_l/\nu_l\),
where \(b_l\) is a signed choice of \(\nu_l/\sqrt{W'_l}\) and \(c\)
is nonzero. Thus the necessary equation is
\[
  3\sum_l b_l Z'_l + I_A\sum_l b_l/W'_l=0.
\]
Write \(M=2+k\) and
\(W(t)=(1+kt,kt-1-k,1-2t)/M\).
Then \(W'=W(y)\), where \(B_k(y)=0\). By Tomar's three-body
dichotomy, the opposite coordinates have one of the following forms.
In the Lagrange case \(Z'=R W(d/y)\), with \(d^3=1\) and \(R^3=M^2\).
This follows from the three distance equations and \(p_{ij}^3=1/M\).
In the Euler case \(Z'=\lambda W(y)\) where
\[
 \lambda^3=F^2,\qquad
 F=2+k\left(\sigma_{13}/y^2-\sigma_{23}/(y-1)^2\right).
\]
Indeed \(p_{ij}=\kappa\sigma_{ij}/W'_{ij}\), \(\kappa^2=1/\lambda\)
and the difference of the first two normalized force equations yields
\(\kappa^3F=1\).

In the Lagrange case set \(U=\sum b_l W_l(d/y)\), \(s=M^2\); in the
Euler case set \(U=\sum b_l W_l(y)\), \(s=F^2\). In both cases put
\(T=\sum b_l/W'_l\).
For the noncollinear equal lift the necessary test is \(3s U^3+T^3=0\);
for the collinear equal lift it is \(U=0\).
In either test the verification may replace \(b_l\) by a signed
\(\nu_l/\sqrt{\ell_l(y;k)}\), removing one common nonzero scalar.

At \(k=3\), evaluate rational interval enclosures for
all 30 roots, all four overall-normalized radical sign choices, each
of three possible \(d\), and each of four possible Euler sign pairs.
Take the rectangular hull of each certified disk and round all arithmetic operations outwards on the grid $2^{-80}$. With $N=2^{80}$ and $q=\lfloor\sqrt{3N^2}\rfloor$, enclose the primitive cube root by real part $-1/2$ and imaginary part $[\lfloor q/2\rfloor,\lfloor(q+2)/2\rfloor]/N$; square this box for its conjugate. For a reciprocal multiply by the complex conjugate and divide by the sum of the two squared real intervals (a squared interval containing zero has lower endpoint zero). A lower bound for the complex modulus is the maximum of the distances of its real and imaginary intervals to zero.
The test includes combinations of shape and triangle type even when those
combinations do not satisfy their corresponding internal equations.
The 840 evaluations in each test give uniform lower bounds
\(2.721333\) for \(\lvert 3sU^3+T^3\rvert\), and \(0.050183\) for
\(\lvert U\rvert\). Thus neither necessary test can hold at \(3\).

For generic parameters, take in each test the product over the four
normalized sign choices. It is invariant under every radical sign
flip, since a global sign changes each odd factor's sign four times.
It is therefore a polynomial in the three radical squares, with rational
function coefficients in \(y,k,d\). Substitute
\(b_l^2=\nu_l^2/\ell_l(y;k)\), clear powers of
\(M y(y-1)\prod_l\ell_l\), and take the product over the three
Lagrange values \(d\), or respectively the four Euler sign pairs.
The resulting numerator polynomials have rational coefficients.
A common denominator may be used throughout. At \(k=3\) no one
of these denominators vanishes on the 30 shape roots: the previous
certificate excludes all \(\ell_l=0\), and \(B_3(0)B_3(1)\ne0\)
by direct evaluation. Their resultants in \(y\) against \(B_k\)
are therefore nonzero at \(k=3\) and vanish at any admissible
specialization exhibiting one of the tests. Multiplying these four
resultants, $P(k)$, and $\Delta(k)$ gives a nonzero obstruction at 3. This proves the proposition.




\subsection{The input to the next reductions}\label{early:input}

Of the 864 surviving marked placements with an assigned core, the
diagram-wide tests remove 88 in diagram 66, 32 in diagram 89, four in
diagram 82 and two in diagram 81. Diagrams 64, 90 and 98 contribute
none of these 864. The 738 remaining marked placements have the following
inventory:
\[
\begin{array}{c|rrrrrrr}
\text{diagram}&20&36&46&47&77&84&93\\
\text{number}&96&20&160&134&38&4&28\\[2pt]
\text{diagram}&94&99&100&101&102&117&\\
\text{number}&72&2&142&28&12&2&
\end{array}
\]
For the separately counted 43 cells without any assigned core, the
original distribution in diagram order $(64,66,88,89,93,94,98)$ is
$(11,11,9,1,1,9,1)$. The whole-diagram exclusions remove 24.
The remaining 19 cells are class 127 cells 0 through 8 (diagram 88),
class 132 cell 53 (diagram 93), and class 147 cells
$5,11,17,23,29,35,41,47,53$ (diagram 94). Each admits six possible
marked slots before further tests. Thus the exact input to the next
step comprises $738+6\cdot19=852$ marked placements.

\paragraph{Finite data underlying this inventory.}
Use revision \texttt{84bbcf2} of the catalogue accompanying~\cite{Tomar}
and its certificate archive (Zenodo record 23472605). Match the cell links in
\texttt{f4/links} to their certificates and apply the Laurent rule to each
right-hand side. For the balanced test use a triple in distinct children of
its least common ancestor in each tree and compare the recorded sets of minima
with the two specified forces. For the clique test compare the recorded minima
with every internal edge and compute the two dimensions above. Remove the
listed whole diagrams and combine with all slots in the 19 specified cells.
This gives the inventory directly from the published cells. The SHA-256 of
the catalogue's decompressed-cell checksum manifest is
\texttt{deb001fef4cf445f53af88d9c06d20be}\allowbreak
\texttt{0ee86dd333a050e50bd9d65fe458b895}; the downstream classifications
can therefore refer to exact triples $(c,i,s)$.
\section{From the specialized catalogue to three diagrams}
\label{br-ledger}
We start with the 738 marked placements with a core and the 114 without a
core obtained above.  A placement means a particular cell and a particular
slot for the mass $t$.  All comparisons in the following two tests are
over $K=\mathbb Q(t)$.  A nonzero minor or obstruction in this field
continues to be nonzero outside the finitely many roots and poles of the
rational functions used.

\paragraph{Linear relations and projected triangles.}
At a node on the path from the origin leaf to the root of one coordinate
tree, give its children coordinates $Y_c$ and set the origin-child
coordinate to zero. Distinct children must have distinct coordinates.
At the root include the row $\sum M_cY_c=0$, where $M_c$ sums all body
masses in child $c$.  The following gives additional rows directly from
the recorded minimizing terms.  Join two bodies when each occurs as a
minimal force in the other's single-body equation in this projection.
On a connected component the minimal orders agree, because opposing
forces differ only by their mass factors. Retain a component if all its
minimal forces are internal. Sum its leading equations with mass
weights. Internal forces cancel. If all the surviving position terms
belong to non-origin children of the node, their weighted sum is a
linear row in its $Y_c$. A nullspace which forces $Y_c=Y_{c'}$ for
distinct children contradicts the tree.

A second test uses three bodies $a,b,s$ in three distinct children at
one node in each tree. Suppose two of their minimal equations in one
projection consist exactly of their two internal forces. The three
internal forces balance there by antisymmetry. If the nullspace
projects with rank one onto their cyclic differences, choose a
representative $(e_a-e_b,e_b-e_s,e_s-e_a)$. Set
\[
 u=m_s^2(e_a-e_b),\quad v=m_a^2(e_b-e_s),\quad
 w=m_b^2(e_s-e_a).
\]
Writing $P_{ab}$ for the leading inverse distance, force balance gives
$P_{ab}^3(e_a-e_b)=\lambda m_s$ cyclically, with the same nonzero
$\lambda$. By the distance equations, the opposite projected
differences are proportional to
$\alpha=P_{ab}/m_s,\ \beta=P_{bs}/m_a,\ \gamma=P_{sa}/m_b$.
They sum to zero; their cubes are proportional to $(1/u,1/v,1/w)$.
Cubing $\alpha^3+\beta^3+\gamma^3=3\alpha\beta\gamma$ gives the test
\begin{equation}\label{br-node}
 (uv+vw+wu)^3=27(uvw)^2 .
\end{equation}
For example, in diagram 47, class 25 cell 57, positions $W_1,W_2,W_3$
are shared, respectively, with bodies $6,5,4$ (one-based labels).
The closed subset $1,2,3$ and the full center of mass give both
$m_1W_1+m_2W_2+m_3W_3=0$ and $m_6W_1+m_5W_2+m_4W_3=0$.
Marking slot 1 or 6 gives $(t-1)W_1=0$ although the tree has $W_1\ne0$.

Here and below, a case count is a finite computation on the catalogue.
For the node tests loop over each input placement, each projection and
each node on the origin path, form the rows just given and reduce them
over $K$. For the remaining placements loop over unordered triples and
both projections and apply \eqref{br-node} when its rank and support
hypotheses hold. The first pass removes 368 core-bearing placements;
the second removes 236 more. The combined removals by diagram are
$20:96$, $36:20$, $46:160$, $47:134$, $77:38$, $100:116$, $101:26$,
$102:12$, and $117:2$. The remaining 134 core-bearing placements are
distributed as follows:
\[
\begin{array}{c|rrrrrrr}
 \text{diagram}&84&93&94&99&100&101&\text{total}\\ \hline
 \text{placements}&4&28&72&2&26&2&134 .
\end{array}
\]
Carry forward all 114 core-free placements. We next remove the
rectangles; all body labels in the rectangle calculations are zero-based.

\subsection{The four rectangle diagrams}
Take $R,S$ to be half separations across the two pairs of rectangle
columns and rows, and write $Z=z/R$, $W=w/S$, $h=(RS)^{-1/2}$ on a
finite Puiseux extension. The normalized forces retain their form in
both coordinates and $p_{ij}^2 Z_{ij}W_{ij}=h^2$. The two different
pairings of the four unit vertices are reciprocal binary balances in
their separating projections. Their sums give opposing coordinates,
so the four vertices limit to $(\pm1,\pm1)$ in all four combinations.
The diagonal products in these coordinates are $4,-4$. Their
contribution at order $h$ to the constant energy is
$K=(\pm1\pm\mathrm i)/2$. Signs need not be physically positive.
The torque identity for any subset $J$ used below follows by
mass-weighted subtraction of its two projected force equations:
\[
 \sum_{j\in J,\ell\notin J}m_jm_\ell p_{j\ell}^3
 (Z_jW_\ell-Z_\ell W_j)=0 .
\]

\paragraph{Diagram 84.}
In class 20, cells 0 and 2, the mixed binary is $0,5$. Its $w$ separation
is $B=w_0-w_5$, with $S/B\to0$, and both binary $z$ coordinates are
$o(R)$. Put $g=(RB)^{-1/2}=o(h)$. The four sides and binary have finite
nonzero inverse distances, the diagonals have order $h$, and the eight
remaining inverse distances have order $g$. For a side separated in
$z$, subtract its endpoint equations there and divide by its separation.
The other finite side at an endpoint has $z$ difference $O(1/S)$,
hence costs $O(h^2)$; diagonal and binary forces cost $O(h^3)$ and
$O(g^3)$. For a side separated in $w$ the external cost is
$O(|g|^3|B/S|)=O(|g h^2|)$. Thus $2p^3=1+O(h^2)$ on every side;
on the mixed binary the same subtraction in $w$ gives
$(1+t)p^3=1+O(g^3)$. Simple cube roots give $o(h)$ changes in all five
finite terms. Subtracting the limit of the energy and dividing by $h$
would give $K=0$. This removes all four placements.

\paragraph{Diagram 100, close column.}
In 24 placements of class 64, one endpoint $u$ of the binary $J$
shares a noncentral $Z$ child with two rectangle vertices $a,b$.
They lie in a common subchild distinct from $u$. At these scales the
binary $W$ coordinates vanish and the corresponding rectangle $W$
coordinates are opposite and nonzero. The two nearest cross edges
$ua,ub$ consequently have leading inverse-distance ratio $k$ with
$k^2=-1$. Every cross-edge torque factor in parentheses has nonzero
limit. Only $ua,ub$ have minimal inverse-distance order; their torque
factors have opposite limits and $m_a=m_b=1$. The subset torque
therefore requires $k^3=1$, an impossibility.

\paragraph{Diagram 100, separated binary.}
In the other two placements (class 64 cell 0) the projected binary
coordinates tend to $(y,0),(-ty,0)$, by adding their weighted binary
equations. The tree gives $y\ne0$ and nonzero entries of
$L=(y-1,y+1,-ty-1,-ty+1)$. Far edges have order $h$. As above the
finite sides satisfy $2p^3=1+O(h^2)$ and the binary satisfies
$(1+t)p^3=1+O(h^3)$ in its separating projection. The order-$h$
coefficient of the energy is
\[
 C(y)/K=1+\sum_{j=1}^4 v_j,\qquad
 v_j^2=4s_jm_j^2/L_j,\quad m=(t,t,1,1),\quad s_j\in\{1,-1\}.
\]
Indeed each pair of edges into a column has sum proportional to
$(\pm1\pm\mathrm i)/\sqrt{L_j}$, and the square of the numerator
divided by $K^2$ is $\pm4$. A cross-edge coefficient $g=\lim p/h$
satisfies $g^2(Z_j-Z_\ell)(-W_\ell)=1$, where $Z_j=y$ or $-ty$.
Differentiation of this explicit local radical gives
$2yg'=g^3Z_j W_\ell$. The torque of the binary divided by $h^3$
therefore gives $C'(y)=0$, while constancy of the energy gives $C(y)=0$.

For fixed $s$ put
\[
 P_s(Y)=(\textstyle\prod_j L_j(Y))^8
 \prod_{\epsilon\in\{-1,1\}^4}\left(1+\sum_j\epsilon_jv_j(Y)\right).
\]
Sign invariance and the separate degree bound eight in each $v_j^2$
show that $P_s\in\mathbb Z[t,Y]$; at infinity its degree is 32 and its
leading coefficient is $t^{16}$. The double zero required above is away
from all denominator zeros. Here is a short alternative to printing the
sixteen discriminants. For each $s$ specialize $t=2$ in
$\mathbb F_{1009}$; take the first 33 integers $Y\ge2$ at which
$D=\prod L_j\ne0$. At each point start from $F(X)=X$ and successively
replace $F$ by $F(X+v_j)F(X-v_j)$ using $v_j^2=4s_jm_j^2/L_j$.
Interpolate $D^8F(1)$ to degree $32$ and take the monic gcd with its
derivative. In lexicographic sign order all sixteen gcds are $1$.
These operations determine the polynomial uniquely by the degree bound
and give a reproducible finite-field certificate that no discriminant
vanishes in $K$.

\paragraph{Diagram 99.}
In class 8 cell 0 the two mixed-binary coordinates vanish at both outer
scales. At least one binary separation in $(Z,W)$ has magnitude
$\gg |h|$ by the distance equation. Subtraction in that projection
therefore gives $(1+t)p^3=1+O(h^2)$, since the far forces are $O(h^3)$.
The four finite sides change by $O(h^2)$ as before. Each binary member
sees four far coefficients containing two choices from $\{1,-1\}$
and two from $\{\mathrm i,-\mathrm i\}$. The energy requires
$K+P+tQ=0$, where $\Re P,\Re Q\in\{-2,0,2\}$. Its real part leaves
only $t=1/4,3/4,5/4$. This removes both generic placements.

\paragraph{Diagram 101.}
For the last two placements (class 116 cell 5), the marked body is 0
or 5. Put $r=m_5/m_0$. The limiting column coordinates of $2,3,5$ are
$1$, of $1,4$ are $-1$, and $Z_0\to-r$ by the $0,5$ binary.
The tree and finite $0,5$ distance give $W_0,W_5=O(h^2)$ in their
common three-child node with the origin. The seven finite edges are
$05,13,24,14,23,25,35$. Subtract the two endpoint equations in $Z$
on $05,13,24$, and in $W$ on $14$: any other finite edge at that end
has coordinate difference $O(h^2)$ in the chosen projection.
Thus those four inverse distances equal their limiting values up to
$O(h^2)$. The summed $W$ equations at $1,4$ give $W_1+W_4=O(h^2)$;
$W_2-W_4,W_3-W_1=O(h^2)$ and the normalization $W_2-W_3=2$ give
$W_2=1+O(h^2), W_3=-1+O(h^2)$. The $W$ equation at 5 now gives
$p_{35}/p_{25}=q+O(h^2)$, $q^3=1$. As $Z_{23}=Z_{25}-Z_{35}$,
\[
 p_{23}^{-2}=(W_2-W_3)\left(\frac{p_{25}^{-2}}{W_2-W_5}
 -\frac{p_{35}^{-2}}{W_3-W_5}\right)
 =2(p_{25}^{-2}+p_{35}^{-2})+O(h^2).
\]
Since $1+q^{-2}\ne0$, $p_{23}^3/p_{25}^3=\gamma+O(h^2)$ for a
nonzero constant $\gamma$. The $W$ equation at 2 is
$1=(2\gamma+m_5)p_{25}^3+O(h^2)$. All seven finite terms therefore
change by $O(h^2)$. The order-$h$ coefficient consists of $K$ and three
terms with respective numerators of the form $(\pm1\pm\mathrm i)$:
their multipliers are $m_5/\sqrt2$, $m_0/\sqrt{1-r}$,
$m_0/\sqrt{-1-r}$. At $(m_0,m_5)=(1,2)$ the sum belongs to
$K+\mathbb Z[\mathrm i]+\sqrt2\,\mathbb Q(\mathrm i)
+\sqrt3\,\mathbb Q(\mathrm i)$; at $(2,1)$ it belongs to
$K+\sqrt2\,\mathbb Q(\mathrm i)+\sqrt6\,\mathbb Q(\mathrm i)$.
In both cases the Gaussian rational component is nonzero, by independent
sign changes of $\sqrt2,\sqrt3$. Taking the product over all radicals
and numerator signs thus gives a nonzero rational obstruction (over
$\mathbb Q(\mathrm i)$), excluding the generic parameter.

\subsection{The fully circled triangle in diagram 93}
\label{br-circled}
There are 28 core-bearing placements of diagram 93 and six more in a
core-free cell. In the following paragraphs we use the labels of
\cite[Theorem 6.15]{Tomar}: $B=\{1,2,3\}$ is the fully circled triple
and $A=\{4,5,6\}$ its complement, with 4 the special vertex.
Chen and Tong derive a signed three-term sum on $B$ for a first limit of a curve through a positive physical point \cite[Prop.~5.16, eq.~(109)]{CT}. For unit masses and projected coordinates proportional to $(1,\omega,\omega^2)$, where $\omega$ is a primitive cube root of unity, the three signed roots may be chosen proportional to $1,\omega,\omega^2$. Their sum vanishes for every $t$ when the marked body lies in $A$.
To see directly what the projected shape and torque conditions on $B$
imply on our mass line, label the triple $(1,1,m)$ and use the shape
variable $x$ of that theorem. This relabeling is permissible because
the three bodies are interchangeable in its hypotheses. Put $d=x-1$,
$M=2+m$, $\ell_1=1+mx$, $\ell_2=\ell_1-M$,
$\ell_3=\ell_1-Mx$, and
\[
 (Y_1,Y_2,Y_3)=(\ell_2\ell_3,\ell_1\ell_3,
 m^2\ell_1\ell_2),\qquad
 \Phi=\textstyle\sum_iY_i^2-2\sum_{i<j}Y_iY_j.
\]
The torque requires $\Phi=0$. The projected three-body shape must
vanish on one of the following eight factors from the same theorem:
\[
\begin{gathered}
 x^2-x+1,\quad x^4-5x^3+9x^2-5x+1,\quad
 x^4+x^3-2x+1,\quad x^4-2x^3+x+1,\\
 2d^3x^2+d^2x^2+ m\epsilon d^3-m\eta x^3-\epsilon d^2-\eta x^2,\quad
 (\epsilon,\eta)=(+,+),(+,-),(-,-),(-,+).
\end{gathered}
\]
If $m=t$, a common zero forces one of eight resultants in $x$ to vanish.
For transparency, at $m=2$ the resultants, taking the shape factor
first, are respectively
\[
\begin{array}{r@{\qquad}r}
1149184 & 2949549457408\\
2949549457408 & 19181142016\\
67997419987533824 & 77003923457572864\\
67997419987533824 & 53648106050289664 .
\end{array}
\]
The entries are read by rows and can be obtained as Sylvester
determinants of the two displayed polynomials. Thus $B$ cannot contain
the marked body at generic $t$. This removes 21 core-bearing and three
core-free placements. If $m=1$, then
$\Phi=9(x^2-x+1)^2$. The centered projected coordinates are proportional
to $(\ell_1,\ell_2,\ell_3)$; for either root of the equilateral
quadratic these are proportional, up to permutation, to
$(1,\omega,\omega^2)$ with $\omega$ a primitive third root of unity.

\paragraph{Five closing equations.}
Rotate reciprocally to $z_5-z_6=1$ and put $h^2=1/w_{12}$ in that
chart. Write $(z_j,w_j)=(Z_j,W_j)$ for $j\in A$,
$(z_l,w_l)=(h^2Z'_l,h^{-2}W'_l)$ for $l\in B$.
The diagram-93 valuation pattern in \cite[proof of Theorem 6.15]{Tomar}
gives bounded capital-letter coordinates and finite nonzero internal
distances; $Z_5,Z_6,W'_l$ have nonzero limits. With $s^2=Z_4$,
$v=h/s$, both $s,v\to0$. For a regular $j=5,6$ the cross inverse
distances are $O(h)$, and for $j=4$ they are $O(v)$ since
$Z_4-h^2Z'_l=s^2(1-v^2Z'_l)$. Denote by $R^Z_j,R^W_j$ the external
forces at $j\in A$. Then
\[
 R^Z_{5,6}=O(h^3),\qquad R^W_{5,6}=O(h),\qquad R^Z_4=O(s^2v^3).
\]
All three $Z$ equations and the two regular $W$ equations consequently
close at the limit. In the inner $Z',W'$ equations the external terms
are also small: in $Z'$ the term from 4 has order
$O(s^2v^3/h^2)=O(v)$ and those from 5,6 have order $O(h)$; in $W'$
they are $O(v^3+h^3)$. Thus the inner triangle solves both isolated
force equations.

Write $a=m_4,b=m_5,c=m_6$ and henceforth use $p_{ij},Z_j,W_j$ for
limits where indicated. The weighted sum of the three $Z$ equations
gives $bZ_5+cZ_6=0$, so $Z_5=c/(b+c)$ and $Z_6=-b/(b+c)$. The
equation at 4 gives $u=p_{45}^3=p_{46}^3$; that at 5 gives
$1=au+(b+c)k$ for $k=p_{56}^3$.
Let $q=p_{46}/p_{45}$, $D_{ij}=p_{ij}^{-2}$ and
$\delta=W_5-W_4$. The distance equations imply $q^3=1$,
$D_{46}/D_{45}=q$, and
\begin{equation}\label{br-outer}
 W_6-W_4=-qc\delta/b,\qquad
 D_{56}/D_{45}=(b+c)(b+qc)/(bc).
\end{equation}
The $W$ equation at 5 now gives $W_4=-k(b-qc^2/b)\delta$ and the
residual at 4 is $(u-k)(b-qc^2/b)\delta$. Here $u\ne k$:
for $q=1$ the distance ratio is at least four; for primitive $q$
its squared modulus is $x^2+x+x^{-1}+x^{-2}\ge4$ with $x=b/c$.
This explicit residual matters because the sixth equation need not
close.

\paragraph{Three further signed placements.}
We include these deletions to give the signed and real ledgers against
the same catalogue. They are dispensable for the energy exclusion
below. If $a_4$ is the order of the original special $z$ coordinate,
then $h,Z_4$ in the rotated chart have orders $2,2+a_4$ respectively.
Pairwise cancellation gives the exact torque
$\sum_{j\in A}m_j(W_jR^Z_j-Z_jR^W_j)=0$. A nonzero limiting $R^W_4$
would force $|s|^2=O(|h|+|s|^2|v|^3)=O(|h|)$ by absorption. Thus for
$a_4<0$ the residual vanishes, which forces $b=c,q=1,W_4\to0$.
In class 132, cells 42, 44 and 51, a regular internal force is minimal
in the special $w$ equation, of order $2$ against the order
$1-3a_4/2$ of an external force. Thus $a_4\le-2/3$. In cell 51 the
position is minimal too, forcing a nonzero limit of $W_4$, a
contradiction.

In cells 42 and 44 the special $W_4$ lies on the origin leaf and hence
vanishes identically. The marked vertex is 4, so $a=t,b=c=1$.
Before taking limits put $x=Z_4$, $S=(Z_5+Z_6)/2$, $d=x-S$,
$D=(W_5-W_6)/2$, $\eta=-(W_5+W_6)/2$; write $D_0=\lim D\ne0$.
The sum of the $Z$ equations gives
$d=(1+t/2)x+o(x)$, since $h=o(x)$. Here
$D_{45}=(1/2-d)(D-\eta)$ and $D_{46}=(1/2+d)(D+\eta)$.
The two limiting signed inverse distances agree as their squares
and cubes agree. Expanding the same local $-3/2$ power yields
$(p_{45}^3-p_{46}^3)/2
=u(3d+3\eta/(2D_0))+o(|d|+|\eta|)$.
The difference at $d=\eta=0$ vanishes for all $D$ near $D_0$.
Set $y=\eta/D_0$, $N=|x|+|y|$. The $Z$ equation at 4 and averaged
$W$ equations at 5,6 give
\[
\begin{pmatrix}
1+u(1+t/2)&3u/2\\ 3tu(1+t/2)&1+tu/2
\end{pmatrix}
\binom{x}{y}=o(N).
\]
For the first row the internal force is
$(p_{45}^3+p_{46}^3)d-(p_{45}^3-p_{46}^3)/2$.
For the second it is
$t[D(p_{45}^3-p_{46}^3)-(p_{45}^3+p_{46}^3)\eta]/2$.
As $x$ has finite order, $N\ne0$ and the limiting determinant $T$
must vanish. By \eqref{br-outer}, $D_{56}/D_{45}=4$, so
$(t+\sigma/4)u=1$, $\sigma=\pm1$. Substitution gives
$(4t+1)^2T=5-36t$ for the positive sign and
$(4t-1)^2T=-3(20t+1)$ for the negative sign. Both fail generically.

We obtain $54,7,126$ signed placements in diagrams $88,93,94$,
respectively (187 total). In diagram 93 the four surviving core-bearing
entries are in cells 43,60,61,62, each marked in catalogue slot 2;
core-free cell 53 retains slots 2,3,4. These are one-based
\emph{catalogue} slots, distinct from the labels used in the force
equations above.

\section{Curves selected by real configurations}
\label{br-real}
Put $M=t+5$, $S=\sum_{i<k}m_i m_k=5(t+2)$, and let $j$ denote the
marked body. On the collision-free algebraic solution set write
\[
 U=\sum_{i<k}m_im_kp_{ik},\qquad I=\sum_i m_i z_iw_i,\qquad
 H=\sum_{i\ne j}p_{ij}+z_jw_j/2 .
\]
Multiply the $z$ force equations by $m_iw_i$ and sum, collecting
the two ends of each edge. The result is
$I=\sum_{i<k}m_im_kp_{ik}^3 z_{ik}w_{ik}=U$. Write $E=I=U$.
Differentiating the distance equation gives
$dp_{ik}=-p_{ik}^3(w_{ik}\,dz_{ik}+z_{ik}\,dw_{ik})/2$.
In $d(U+I/2)$ the coefficients of every $dz_i,dw_i$ cancel by the
force equations. Only the $t$ derivatives of the mass factors remain.
Thus, on the full mass-parametrized set,
\begin{equation}\label{br-env}
 \tfrac32\,dE=H\,dt .
\end{equation}

\paragraph{Specialization and choice of curve.}
Work over $k=\overline{\mathbb Q}$ in the normalized algebraic slice,
which is of finite type. On an irreducible component dominating the
$t$-line, the image of $dE$ in the relative differentials of its
function field over $k(t)$ vanishes by \eqref{br-env}.
Characteristic zero makes $E$ algebraic over $k(t)$: a
transcendental $E$ could be extended to a separating transcendence
basis and the derivation taking $E$ to $1$ would contradict $dE=0$.
The algebraic derivative $E'$ satisfies $H=3E'/2$ in the function
field. For example, for the minimal polynomial $P(t,E)=0$ it is
$-P_t/P_E$, with $P_E\ne0$ by separability. Clearing denominators in
the polynomial equations for $E,H$ gives identities on the entire
component. Discard vertical components' parameter values and the
finitely many specializations at which an equation becomes the zero
polynomial. For each other fixed $t$, each quantity takes finitely
many values even at singular points. Reciprocal rotations preserve
both $E,H$. Now fix a positive real $t_0$ outside the omitted algebraic values, base-change the slice to $\mathbb C$, and take closures in the fibre at $t_0$. Passing from a physical configuration to the algebraic rotation slice need not preserve $w_i=\overline{z_i}$, but it preserves the zero coordinates in each projection and the two scalar values.

At a real centered configuration $z_i=0$ if and only if $w_i=0$.
At most one body is at the origin. Partition any infinite family
of real shapes at fixed $t$ by that body's identity (or its absence)
and by the values of $E,H$. Some stratum is infinite.
Take its relative Zariski closure inside the locus prescribing both
zeros and $z_iw_i\ne0$ elsewhere. A positive-dimensional component
has constant $E,H$ with their values attained at real configurations.
Some squared distance varies there, by the finite-fibre property of
the distance map on the rotation slice. Intersect through a smooth
point where that distance has nonzero differential with general
hyperplanes to get an irreducible curve still varying in that
distance. This is a curve to which the nonpositive-place argument
applies. The prescribed zeros are identities of the curve; a
nonzero rational function may have positive finite order but occupies
a distinct leaf from the origin. Hence the two origin leaves of
the place contain exactly the same set of bodies. Complex conjugacy
at the place is unnecessary.

For clarity, applying this last test to the three signed lists and
identifying marked placements under permutations preserving the
colored graph, optionally exchanging its projections, gives
\begin{center}
\begin{tabular}{lrrrr}
\hline
 &88&93&94&total\\\hline
signed placements&54&7&126&187\\
signed classes&8&6&41&55\\
matching origin leaves, placements&30&6&54&90\\
matching origin leaves, classes&5&5&19&29\\
\hline
\end{tabular}
\end{center}
For this quotient, transport both rooted trees (including the origin),
the marked body and every support of minimizing terms; compare the
result with the recorded target cell. This finite permutation check
uses at most $2\cdot6!$ images per placement.

\paragraph{Bounds at physical points.}
In a real configuration write $r_{ik}=|q_i-q_k|$. Centering gives
$I=M^{-1}\sum_{i<k}m_im_k r_{ik}^2$. Jensen's inequality for
$x^{-1/2}$ with the weights $m_im_k/S$ gives
$E\ge S\sqrt{S/(ME)}$, hence
\begin{equation}\label{br-jensen}
 E\ge 5(t+2)/(t+5)^{1/3}.
\end{equation}
This is also the direct specialization of the positive-configuration bound $M^{1/3}E\ge\sum_{i<k}m_i m_k$ used by Chen and Tong \cite[proof of Prop.~5.21, following (112)]{CT}: the total mass is $5+t$ and the ten unit pairs and five marked pairs contribute $10+5t$.
Also $t\sum_{i\ne j}p_{ij}\le E$. Since
$-tq_j=\sum_{i\ne j}q_i$, Cauchy's inequality gives
$t^2|q_j|^2\le5(I-t|q_j|^2)$. We obtain
\begin{equation}\label{br-elastic}
 0<\frac{2tH}{3E}\le\frac{2t+15}{3t+15}.
\end{equation}
Both inequalities hold for the values of the chosen curve at any
place where they are evaluated as its constant invariants.

We will also differentiate candidate end energies. For each of the
finitely many algebraic candidates $e(t)$, take a finite cover of
an open parameter line on which its value is a regular function.
Impose $E=e$ on the pulled-back solution set. On every dominating
component of this incidence, differentiating this equality and using
\eqref{br-env} gives $2H/3=e'$ identically. Discard the finitely many
vertical images, branch values and poles and clear denominators.
It follows that $2H/3=e'$ at every remaining point of the incidence.
This argument applies even if a putative real family is known only
at one fixed $t$, since $E$ is constant along the curve and its value $e(t)$ at the end therefore places the points of the curve
in the algebraic incidence. Separate finiteness of the vertical images
is essential to the cofinite assertion.

\subsection{Energy exclusion of diagram 93}
The completely circled triple of Section~\ref{br-circled} has unit
masses. Its limiting energy $B_0$ takes values in a finite
$t$-independent set and
\begin{equation}\label{br-boundB}|B_0|\le3^{2/3}.\end{equation}
Indeed its two coordinate vectors solve the isolated equations and
are centered. If they are independent, the force operator
$L(v)_i=\sum_{k\ne i}p_{ik}^3(v_i-v_k)$ acts as the identity on
the entire centered two-plane and annihilates the constant vector.
Thus its off-diagonal entries are $-1/3$, all $p_{ik}^3=1/3$,
and there are finitely many signed choices. If the vectors are
proportional, the bilinear inertia vanishes, since one of them
is proportional to $(1,\omega,\omega^2)$ and
$1+\omega^2+\omega^4=0$. The isolated identity $I=U$ then gives $B_0=0$.

For the complementary triple set $g=p_{56}/p_{45}$.
The five closing equations and \eqref{br-outer} give
\begin{equation}\label{br-g}
 g^2=\frac{bc}{(b+c)(b+qc)},\qquad
 p_{45}^{-3}=a+(b+c)g^3.
\end{equation}
The cross inverse distances vanish, so $E=A+B_0$. Using \eqref{br-g},
the complementary triple's energy is
\begin{align}
 A&=p_{45}\{a(b+qc)+bcg\}=(b+qc)p_{45}^{-2},\label{br-A}\\
 |A|^{3/2}&\le a|b+qc|^{3/2}+(bc)^{3/2}/\sqrt{b+c}.\label{br-boundA}
\end{align}
We can now eliminate all placements, regardless of whether the
sixth limiting force equation closes.

\paragraph{Mark at the special vertex.}
With $a=t,b=c=1$, \eqref{br-boundA} yields
$|A|\le2(t+1/4)^{2/3}$. By weighted arithmetic--geometric mean,
\[
 (t+5)^{1/3}(|A|+|B_0|)
 \le 2\frac{(t+5)+2(t+1/4)}3+\frac{(t+5)+2\cdot3}3
 =\frac{7t+22}{3}<5(t+2),
\]
contradicting \eqref{br-jensen}.

\paragraph{Mark at a regular vertex.}
Assume $b=t,a=c=1$ after relabeling. Equation \eqref{br-boundA} gives
\[
 |A|\le(t+1)\left(1+\frac{t^{3/2}}{(t+1)^2}\right)^{2/3}
 \le\tfrac54(t+1).
\]
The fraction has maximum $3\sqrt3/16<1/3$ at $t=3$.
Use $|B_0|<21/10$. On $0<t\le2$ the cube root $(t+5)^{1/3}$ is
less than $2$; multiplying this energy estimate by $2$ gives
a strict upper bound $5(t+2)$. On $2\le t\le32$ use
$(t+5)^{1/3}<10/3$ and the same upper bound follows (at $t=2$
the linear inequality is strict). Both intervals contradict
\eqref{br-jensen}.

For $t\ge32$ let $k=(t+1)g^3$, so
$k^2=t^3/[(t+1)(t+q)^3]$. By \eqref{br-A},
$A=(t+q)(1+k)^{2/3}$ with a locally chosen cube root.
Taking its logarithmic derivative gives
\[
 1-\frac{t A'}{A}=\frac{q/(t+q)+kt/[3(t+1)]}{1+k}.
\]
Here $|t+q|\ge31$, $|k|\le(32/31)^{3/2}/\sqrt{33}<1/5$, whence
\[
 \left|1-\frac{tA'}{A}\right|
 \le\frac{1/31+1/15}{4/5}<1/8,\qquad |A|\ge31(4/5)^{2/3}>24 .
\]
Thus $|B_0|<|A|/10$. There are finitely many choices of $q$, the
radicals and $B_0$, so the incidence argument applies to all of them.
At the remaining parameters $2H/3=A'$. Consequently
$|tA'|/|E|>(7/8)/(11/10)=35/44>3/4$, whereas
\eqref{br-elastic} gives $2tH/(3E)<3/4$ on this interval.
This eliminates all six physical placements of diagram 93.

\subsection{The marked binary of diagram 88}
\label{br-markbinary}
We prove the following lower bound at every real configuration:
\begin{equation}\label{br-Hb}
 H>\frac{3}{2(1+t)^{1/3}} .
\end{equation}
Choose a pair of orthogonal real axes avoiding all finitely many
projected coincidences. In one projection put
$a_{il}=(|x_i-x_l|/r_{il})^3\le1$ and hold these numbers fixed.
In the chamber of the order of the $x_i$ the function
\[
 V(u)=\sum_{i<l}\frac{m_im_la_{il}}{|u_i-u_l|}
       +\tfrac12\sum_i m_i u_i^2
\]
is strictly convex: each inverse separation is convex on the chamber
and the moment term is strictly convex. Its gradient at $x$ vanishes
by the projected force equations. Thus $V(x)$ is its minimum, and
differentiating along $s x$ at $s=1$ gives
$V(x)=3I_x/2$, $I_x=\sum_i m_i x_i^2$.

For an upper bound on the minimum, use the trial vector
$u_i=s(\operatorname{rank}(i)-\operatorname{rank}(j))$, $s>0$.
Let $B$ sum $m_im_l$ divided by the rank gap over all pairs and let $C$
sum the squares of the five rank differences from $j$. Bounding
$a_{il}$ by 1 gives $V(u)\le B/s+Cs^2/2$.
Minimization gives $I_x\le(B^2C)^{1/3}$. Deleting the marked rank
can only shorten the other rank gaps, so
\[
 B\le4+3/2+2/3+1/4+5t\le(77/12)(1+t),\qquad
 C\le1^2+2^2+3^2+4^2+5^2=55 .
\]
The bound for $C$ follows, for example, by ordering the five absolute
rank differences, whose $k$th smallest is at most $k$.
Adding the two projection estimates gives
$I\le2(55(77/12)^2)^{1/3}(1+t)^{2/3}<27(1+t)^{2/3}$.

Centering gives $\sum_{i\ne j}r_{ij}^2=I+(t+5)|q_j|^2$.
Put $y=|q_j|^2(1+t)^{1/3}$. If $y\ge3$, the marked position term
in $H$ and positivity of the inverse distances prove \eqref{br-Hb}.
If $y<3$, Jensen's inequality with five equal weights and
$(t+5)/(t+1)\le5$ give
\[
 (1+t)^{1/3}\sum_{i\ne j}1/r_{ij}
 \ge\frac{5\sqrt5}{\sqrt{27+5y}}>3/2 ,
\]
again proving the bound.

In diagram 88 the binary closes in both projections
\cite[Lemma 6.16(a) and its proof]{Tomar}. If its masses are $t,1$,
subtracting the two limiting binary equations gives the common
inverse distance $b^3=(1+t)^{-1}$; their mass-weighted sum gives
center zero at each leading scale. Therefore
$\lim z_jw_j=b^{-2}/(1+t)^2=b/(1+t)$ by the distance equation.
The four external inverse distances at $j$ vanish. The constant
value $H$ must satisfy
\[
 |H|=|b|\left(1+\frac1{2(1+t)}\right)
 \le\frac{3}{2(1+t)^{1/3}},
\]
contrary to \eqref{br-Hb}. The marked binary accounts for ten
placements in two symmetry classes. After these and diagram 93
have been removed, 74 placements remain: 20 in diagram 88 and
54 in diagram 94, representing three and nineteen symmetry classes,
respectively. We now address the geometry common to these remaining
classes.
\section{The signed four-body system of diagram 94}\label{nf:signed}

Denote the shared vertex by $c$, the outer vertices of its four-body cluster by
$1,2,3$, and the two vertices of the opposing cluster by $5,6$.
A permutation of the three outer vertices or of the two binary vertices preserves
the diagram. \begin{proposition}\label{nf:prop}
Outside a finite set of positive values of $t$, the mass tuple consisting of
five unit masses and $t$ has no normalized complex singular place of diagram
94, in any ordering or with any compatible choice of distance signs.
\end{proposition}
The unequal mass has three possible roles. For a generic mass at
$c$ or in the binary, the leading equations admit two particular geometries. Excluding their lifts requires the six-body equations. For an outer unequal
mass we can already derive a contradiction from the four-body equations at
$t=\infty$.

We use the necessary leading equations for this diagram from
\cite[\S6.7, Lemmas 6.16(b) and 6.17(b)]{Tomar}.
For a first limit on a curve through a positive physical point, Chen and Tong obtain a necessary ambient mass condition by separating the possibilities for the three inverse-distance cubes incident to the shared body \cite[Prop.~5.20]{CT}. Here we keep the signs and specialize the quartet equations and their lifts on the five-equal-mass line.
In the large projection take $X_c=0$, and normalize one outer $X$ to $1$.
Write $X_{ij}=X_i-X_j$, $Y_{ij}=Y_i-Y_j$.
All six internal $X_{ij},Y_{ij}$ are nonzero. Each edge has its own signed inverse
distance $p_{ij}=p_{ji}$ with $p_{ij}^2X_{ij}Y_{ij}=1$.
The quartet equations used below are
\begin{align}
 X_i&=\sum_{j\ne i}m_j p_{ij}^3 X_{ij} && (i=1,2,3,c),\label{nf:forcex}\\
 Y_i&=\sum_{j\ne i}m_j p_{ij}^3 Y_{ij} && (i=1,2,3).\label{nf:forcey}
\end{align}
Here the sums run only through the quartet. There is no need to impose a
fourth small-coordinate equation. The weighted sum of \eqref{nf:forcex}
gives $\sum_{i=1}^3m_iX_i=0$. The distant edges give a signed torque
\begin{equation}\label{nf:torque}
 \sum_{j=1}^3 \frac{m_j}{\sqrt{X_j}}
 (\epsilon_j+\mathrm i\rho\eta_j)=0,\qquad
 \rho=(m_6/m_5)^{3/2},\qquad \epsilon_j,\eta_j\in\{1,-1\}.
\end{equation}
A choice of square root of each $X_j$ is absorbed in the two signs.
For positive masses the real positive value of $\rho$ is sufficient since
the remaining sign is already allowed.

\subsection{The unequal mass at the shared vertex}\label{nf:center}
Here $m_c=t$. The three outer and two binary masses are $1$.
With the two equal binary masses, the torque equations and center relation give just fourteen possible labelled triples $X$. Indeed write the torque terms as $v_i=d_i/\sqrt{X_i}$, with $d_i=\pm(1\pm\mathrm i)$, put $z=v_2/v_1$ and $e_i=d_i^2/d_1^2\in\{1,-1\}$. Then
\[
z^2(1+z)^2+e_2(1+z)^2+e_3 z^2=0,\qquad
X=(1,e_2/z^2,e_3/(1+z)^2).
\]
The $(+,+)$ polynomial is $(z^2+z+1)^2$. The other three quartics each have four distinct roots. Modulo 673 these polynomials split completely, with fourteen distinct admissible projected triples. This verifies the corresponding distinctness and count in characteristic zero. Every projection is independent of $t$.

Put $p_*=p_{1c}$, $r_i=p_{ic}/p_*$, $q_{ij}=p_{ij}/p_*$, $v=p_*^{-3}$, $y_i=p_*^2(Y_i-Y_c)$, and $c_0=p_*^2Y_c$. The distances, the center $X$ equation, and the outer $X$ and $Y$ equations give
\[
r_1=1,\quad \sum X_i r_i^3=0,\quad y_i=(X_i r_i^2)^{-1},\quad
q_{ij}^2=[(X_i-X_j)(y_i-y_j)]^{-1},
\]
\[
a_i=\sum_{j\ne i}q_{ij}^3(X_i-X_j)/X_i,\quad
v=t r_i^3+a_i,\quad
C_i=\sum_{j\ne i}q_{ij}^3(X_jy_i/X_i-y_j)=v c_0 .
\]
To see the last identity, subtract $y_i$ times $v=t r_i^3+a_i$ from the outer equation $v(y_i+c_0)=t r_i^3y_i+\sum q_{ij}^3(y_i-y_j)$. If some $r_i^3\ne1$, subtracting the outer $X$ equations determines $t=(a_i-a_1)/(1-r_i^3)$ uniquely at that geometry.

For each prescribed $X$, the diagonal cubic $\sum X_i r_i^3=0$ in projective ratio coordinates is smooth and geometrically irreducible, since all its coefficients are nonzero. Restrict its $r_1=1$ chart to the open set with nonzero ratios and collision-free $y$'s. Adjoining the three $q_{ij}$ gives a finite étale cover there, so each irreducible component maps onto this open curve. To show the simultaneous equations $C_1=C_2=C_3$ cut out only finitely many points, a single base point whose entire fibre violates at least one equality suffices for each $X$.

Here are explicit witnesses modulo $673$. Each column gives $(X_2,r_2,r_3)$ with $X_3=-1-X_2$ and $r_1=1$.
\[
\begin{array}{r|rrrrrrr}
X_2&6&53&96&111&127&162&255\\ r_2&2&5&4&3&6&3&4\\ r_3&138&73&10&107&329&4&229\\ \hline
X_2&417&510&545&561&576&619&666\\ r_2&5&2&5&3&2&13&5\\ r_3&178&167&60&211&250&18&353
\end{array}
\]
For each column evaluate $y_i$ in the prime field, take both roots of
each $q_{ij}^2$ in $\mathbb F_{673^2}$ and evaluate the three $C_i$.
In all fourteen columns the count of sign vectors with common $C_i$ is
zero. A common reversal of the three signs simply negates all the $C_i$. For the transfer, choose a place over 673 for each algebraic shape. The modular value of $r_2$ lifts as an ordinary nonzero integer, and the cubic for $r_3$ has the tested simple nonzero residual root. Thus it has a lift over a finite extension of the 673-adic numbers. Every denominator at this point is a unit. The reductions of its square-root sheets occur among the tested sheets. Hence no lift of the test point has both required equalities. A component supporting both equalities identically would have to meet that fibre. Consequently there are finitely many sign-and-ratio tuples with the required equalities for each $X$. Outside $r_i^3=1$ for every $i$, this also makes $t$ finite.

\paragraph{The two persistent patterns.}

If all $r_i^3=1$, the force scales allow variable $t$ only if $a_1=a_2=a_3$. The parameters $r_2,r_3$ now range over nine ordered pairs. For each of the fourteen columns, substitute each of the nine ordered pairs of cube roots of unity, omit colliding ordinates, and run the same eight-sign loop retaining sheets with both common $a_i$ and common $C_i$. The only survivors are $X=(1,\omega,\omega^2)$ in either orientation, with every $q_{ij}^3$ equal, and
\[
(r_1,r_2,r_3)=(1,1,1),\quad y_i=X_i^{-1},
\qquad\text{or}\qquad
r_i=X_i^{-1},\quad y_i=X_i .
\]
Reduction is injective on the fourteen possible projected triples and on the cube roots of unity. On any fixed collision-free tuple, the two signs of each nonzero square root also remain distinct. In the exceptional tuples, a collision $y_i=y_j$ modulo 673 requires $X_i/X_j$ to be a cube root of unity modulo 673. Since $X_1+X_2+X_3=0$ and the three coordinates are distinct, this implies the equilateral projection there. That projection lifts uniquely as a projected triple, all its coordinates are cube roots of unity, and the residual collision is a genuine collision in characteristic zero. For both surviving patterns the three squared edge cube values are equal in characteristic zero, so equality of their cubes modulo the odd prime forces the same signs there and in characteristic zero. In either pattern $C_i=0$, hence $Y_c=0$.

The common outer cube $\kappa$ equals $\pm1/(3\sqrt3)$ in the first
pattern and satisfies $\kappa^2=-1/27$ in the second. These identities follow
by cubing the squared-distance equations. In each case $v=t+3\kappa$,
since $\sum_jX_j=0$. Section~\ref{nf:lift} treats both patterns.

\subsection{The unequal mass in the opposing binary}\label{nf:binary}
Here the quartet has four unit masses. It is useful first to prove that the
set of possible large projections is finite without using the torque equation.
Write $X_c=0$, $X_1=1$, and $X=(1,x,-1-x)$ for the three outer abscissas. Their sum vanishes by the sum of the four abscissa equations. All $X_i$ and $X_i-X_j$ are nonzero. Put $p=p_{1c}$, $r_i=p_{ic}/p$, $s_i=r_i^3$, $r=(1,v,w)$, $K=p^{-3}$, and
\[
 U_i=p^2(Y_i-Y_c)=1/(X_i r_i^2),\qquad
 E_{ij}=(p_{ij}/p)^3,\qquad E_{ij}^2=D_{ij}^{-3},\quad
 D_{ij}=(X_i-X_j)(U_i-U_j).
\]
The center abscissa equation gives $\sum s_i X_i=0$. The outer equations give, with a common constant $c$,
\[
 \sum_{j\ne i}E_{ij}(X_i-X_j)=(K-s_i)X_i,\qquad
 \sum_{j\ne i}E_{ij}(U_i-U_j)=(K-s_i)U_i+c.
\]
Put $P_i=(X_i,U_i)$. For an affinely independent outer triple, let $\lambda_i$ be the barycentric coordinates of $(0,0)$ relative to $P_i$. With $\bar P=\sum P_i/3$, summing the force equations lets us write their two-component right-hand sides as
\[
 K(P_i-\bar P)-s_iP_i+\tfrac13\sum_j s_jP_j .
\]
Let $S=\sum s_i/3$. The coefficients in this formula sum to $S-s_i$. Subtracting $(S-s_i)\lambda_j$ from each coefficient leaves the vector unchanged and makes their sum zero. Uniqueness of this representation for an affinely independent triple gives, for $i\ne j$,
\[
 E_{ij}=(K-s_j)/3+(S-s_i)\lambda_j.
\]
The centered vectors $(\lambda_i-1/3)_i$ and $(s_i-S)_i$ both have sum zero and pair to zero with $X$. This is a one-dimensional subspace because the three abscissas are distinct. Taking the corresponding vanishing minors now gives
\[
 E_{ij}-E_{ik}=-\lambda_i(s_j-s_k)\quad (i,j,k\text{ distinct}).
\]
For cyclic $(i,j,k)$, put $n_i=X_jU_k-X_kU_j$. Then $\lambda_i=n_i/\sum n_j$.

If $v^3\ne w^3$, the center equation gives $x=(w^3-1)/(v^3-w^3)$. Substitute it in $D_{ij}$ and $\lambda_i$. Put $P(A,B,d)=(A-B)^2-2(A+B)d+d^2$. It vanishes if $A=e^2$, $B=f^2$, and $d=(e-f)^2$. Consequently the following two expressions vanish:
\[
 P(D_{12}^{-3},D_{13}^{-3},\lambda_1^2(v^3-w^3)^2),\qquad
 P(D_{12}^{-3},D_{23}^{-3},\lambda_2^2(1-w^3)^2).
\]
To clear each expression, write its three inputs as factored rational functions and divide them all by their common factor with exponent equal to the minimum of each irreducible factor's three exponents. Every removed numerator factor divides the first input's numerator. All removed denominator factors occur in the denominators of the original chart coordinates, distances or barycentric coordinates. These factors are units on this chart. The resulting bivariate integer polynomials have bidegrees $(102,102)$ and $(96,102)$, respectively, with 6827 terms each. Their exact polynomial gcd is 1. They therefore have finitely many common zeros over the algebraic closure, and the projected ratios in this chart form a finite set.

If $v^3=w^3$, the center equation makes both cubes 1. For a noncollinear triple the difference identity then makes all $E_{ij}$ equal. For each fixed choice of cube roots $r_i$, the resulting identity $D_{12}^3=D_{13}^3$ is nonidentical in $x$. Indeed at $x=1$, the first distance factor vanishes while $X_3=-2$ and $U_3=-r_3/2\ne U_1=1$. Hence this chart has finitely many projections as well.

If the outer points are affinely collinear, write $U_i=A X_i+B$. Subtract $A$ times the abscissa equations from the ordinate equations and take row differences. This gives $B(s_i-s_j)=0$. If $B=0$, $A=1$, so $r_i=\sigma_i/X_i$ with $\sigma_i=\pm1$. Now $\sum_i\sigma_i/X_i^2=0$ by the center equation, a nonidentical rational condition with a pole at $x=0$. If $B\ne0$, all $s_i=1$. For each fixed choice of $r_i$, collinearity of $(X_i,r_i/X_i)$ gives a nonidentical rational condition because as $x\to0$ the second ordinate has a simple pole whose coefficient in the collinearity determinant is multiplied by $X_1-X_3\to2$. This proves finite projection in every case.

\paragraph{The binary torque selects a regular projection.}

For each of the finitely many projections and square-root choices, the torque equation has the form $A+\mathrm i\rho B=0$, where $A,B$ are signed sums of $1/\sqrt{X_i}$ and $\rho=(m_6/m_5)^{3/2}$. Since one binary mass varies, this can hold for infinitely many positive $t$ only if $A=B=0$. From the signed sum $A=0$, normalize its terms as $(1,z,-1-z)$. Their inverse squares are $X_1,X_2,X_3$. Multiplying $\sum X_i=0$ by $z^2(1+z)^2$ gives $(z^2+z+1)^2=0$. Thus, up to permutation, $X=(1,a,b)$ with $a^2+a+1=0$ and $b=a^2$.

\paragraph{Center edges on the regular projection.}

If the outer triple is noncollinear, the center equation now gives $w^3=-a-bv^3$. In the ring $\mathbb Z[v,a,w]/(a^2+a+1,w^3+a+bv^3)$ the three barycentric minors, multiplied by $v^2w^2$, are
\[
 n_1=bv^2-aw^2,\quad n_2=bv^2w^2-av^2,\quad n_3=bw^2-av^2w^2 .
\]
For the edges $(12,13,23)$, write the corresponding $D_i=N_i/T_i$. The three numerators are
\[
 (1-a)(v^2-b),\quad (1-b)(w^2-a),\quad (a-b)(bw^2-av^2),
\]
and the denominators are $v^2,w^2,v^2w^2$. Let $n=\sum n_i$. Use the pairs $(i,j)=(12,13)$ and $(12,23)$ with $(k,F)=(1,v^3-w^3)$ and $(2,1-w^3)$, respectively. Their polynomial relations are
\[
 P\left(T_i^3N_j^3n^2,\ T_j^3N_i^3n^2,\ F^2 n_k^2 N_i^3N_j^3\right)=0 .
\]
For the exact integer calculation put $b=-a-1$. Reduce the three arguments and then $P$ modulo the monic $w^3+a+bv^3$ and $a^2+a+1$, in that order. Take the resultant with $w^3+a+bv^3$ in $w$, reduce modulo $a^2+a+1$, and take the resultant with $a^2+a+1$ in $a$. These two integer polynomials in $v$ have degrees 264 and 312, and their primitive gcd is exactly $(v-1)^4(v^2+v+1)^{21}$. Every original solution must annihilate both, giving $v^3=w^3=1$.

For collinear triples with $B\ne0$ this was proved above. If $B=0$, the center equation gives $0=\sum\sigma_i/X_i^2=\sum\sigma_i X_i$. The sole rational relation among $1,a,b$ is proportional to their sum, so all signs are equal. Since $r_1=1$, every cube is again 1.

With $s_i=1$ the large-coordinate equations force equality of all three
$a_i$ of Section~\ref{nf:center}. The remaining ordinate equations impose
the same $C_1=C_2=C_3$. The finite sign enumeration in that section applies
because it uses only $X,r,q_{ij}$, all internal to this quartet. Thus both
mass placements reduce to the following lifting problem.

\section{Why neither persistent diagram-94 geometry lifts}\label{nf:lift}
In each of the two patterns the large outer projection is
$(1,\omega,\omega^2)$, $\omega^2+\omega+1=0$. In the complexified
coordinates the small projection is proportional to $(1,\omega^2,\omega)$
in the first pattern and to $(1,\omega,\omega^2)$ in the second. We call
them the equilateral and collinear patterns, respectively. The common scale
and all signs of the inverse distances remain algebraic variables.
Both patterns satisfy the leading equations and the first torque condition.
To distinguish a formal leading configuration from an actual singular branch,
we calculate the third-order terms of a reduced stationary potential.

\subsection{Coordinates and the distant-edge expansion}
Write $s=4$, $A=\{5,6\}$ and $B=\{1,2,3\}$, $S=m_s$, $m=m_5+m_6$, $n=3$, and $M=m+n+S$. In coordinates relative to $s$, write the positions as

\[
(z_i-z_s,w_i-w_s)=(hU_i,h^{-1}(v_i+h^2b)),\quad i\in A,
\]
\[
(z_j-z_s,w_j-w_s)=(h^{-1}(x_j+h^2a),hY_j),\quad j\in B.
\]

Choose $\sum_A m_i v_i=\sum_B x_j=0$, $v_5-v_6=1$ and $x_1=1$. The two difference scales can be chosen after a common complex rotation and a finite square-root extension of the local parameter. Both the leading binary separation in $w$ and the large $z$-projection of the first outer body relative to its cluster mean are nonzero. Their product determines $h^{-2}$. In particular $v_5=m_6/m$, $v_6=-m_5/m$, and $x_3=-1-x_2$. The orders from diagram 94 imply $h\to0$ and bounded small coordinates $U_i,Y_j$. The actual centered coordinates of $s$ are $O(h)$ in both projections. Summing all six absolute positions with the mass weights shows that their relative large-cluster means are $O(h^2)$ in the displayed scaled coordinates, so $a,b$ are bounded too. Distinct internal edges have nonzero limiting squared distances. The nonzero products give analytic continuations of both square roots on each internal edge near the seed.

Write $I$ for the centered mass-weighted sum of the products of the two absolute coordinates. Subtracting the square-product of the relative center of mass gives the exact formula
\[
I=\sum_A m_i U_i v_i+\sum_B x_jY_j+
 h^2\left[bU_A+aY_B-\frac{(U_A+na)(mb+Y_B)}{M}\right],
\quad U_A=\sum_A m_iU_i,\quad Y_B=\sum_B Y_j.
\]
For the signed potential $V=\sum_{i<j}m_i m_jp_{ij}$, the force equations with unit multiplier give stationarity of $L=V+I/2$. Indeed on a fixed square-root branch one has $\partial_{w_i}p_{ij}=-p_{ij}^3(z_i-z_j)/2$. Thus $\partial_{w_i}L=0$ is the $z$-force equation, and $\partial_{z_i}L=0$ is its $w$ analogue. The pullback of $L$ to the chart is therefore stationary at every genuine solution with $h\ne0$, including in the $a,b,h$ directions.

At fixed chart coordinates expand $L=F+hG+h^2 H+h^3 J+O(h^4)$. For clarity the constant term is
\begin{align*}
 F={}& S\sum_{i\in A}m_i p_{si}+m_5m_6p_{56}
      +S\sum_{j\in B}p_{sj}+\sum_{j<k\in B}p_{jk}\\
    &+\tfrac12\left(\sum_{i\in A}m_iU_iv_i+\sum_{j\in B}x_jY_j\right),
\end{align*}
with the internal inverse distances evaluated at $h=0$. Write $g_{ij}^2=-1/(x_jv_i)$ with the signed branches determined by the distant edges. Then
\[
G=\sum_{i\in A,j\in B}m_i g_{ij},\qquad
J=\frac12\sum_{i,j}m_i g_{ij}
\left(\frac{U_i-a}{x_j}+\frac{Y_j-b}{v_i}\right).
\]
These follow by scaling the exact squared distances on cross edges. All internal distances depend only on $h^2$ at fixed chart coordinates. If $D_A=\sum_A m_i p_{si}/v_i$ and $D_B=\sum_B p_{sj}/x_j$ denote their internal expressions evaluated at $h=0$, then
\[
2H=b(U_A-SD_A)+a(Y_B-SD_B)-(U_A+na)(mb+Y_B)/M.
\]

\subsection{Elimination of seven stationary variables}

Put $x=(1,\omega,\omega^2)$. At $h=0$ the outer small coordinates are
$Y_j=\beta y_j$, with $y_j=1/x_j$ in the equilateral case and $y_j=x_j$
in the collinear case. The center-edge ratios $r_j$ are respectively $1$ and
$1/x_j$. Write $p_*=p_{s1}$. The quartet equations
give $p_*^{-3}=S+3\kappa$ and $\beta=p_*^{-2}$ in both cases. In particular,
$Y_B=D_B=0$. The latter follows from $D_B=p_*\sum r_j/x_j=0$.
The large and small quartet force equations give the stationary equations of
$F$ in the three free $Y$ variables and the one free $x$ variable.
The two $U$ equations impose the binary stationarity conditions below.
At the common regular projection, $g_{ij}=d_i\sigma_{ij}x_j$, where $d_i^2=-1/v_i$ and $\sigma_{ij}=\pm1$. A nonzero signed sum of $1,\omega,\omega^2$ is twice a sixth root of unity. The ratio of two such sums is again a sixth root of unity, and none is purely imaginary. For positive binary masses $m_6d_6/(m_5d_5)$ is purely imaginary. Hence $G=0$, the necessary $h$-stationarity condition at zero, implies \emph{each} row sum vanishes, with constant signs in each row. Absorb these signs in $d_i$. For infinitesimal centered shape changes one has $\sum_j dx_j=0$, and $dg_{ij}=-(d_i/2)dx_j$ at the seed. Thus $G=0$ and $dG=0$ there.

Solve $\partial_aL/h^2=\partial_bL/h^2=0$ for $a,b$ by the analytic implicit function theorem. At $h=0$ the diagonal of their Jacobian is zero and its off-diagonal is $-mn/(2M)\ne0$. At the seed, the unique solution is
\[
b_0=0,\quad \lambda=(U_A-SD_A)/m,\quad a_0=(M\lambda-U_A)/n.
\]
For fixed $x_2$, the Hessians of $F$ in the two $U_i$ and the three $Y_j$ are invertible at the all-equal witness, as checked below. This holds outside finitely many $t$ on each algebraic family of these seeds. Eliminate these five variables by their stationary equations. At fixed $x_2$ their first derivatives in $h$ vanish at zero because $G$ depends only on $x$. Consequently the reduced coefficients of $h^2$ and $h^3$ are $H$ and $J$ evaluated at the unperturbed transverse stationary values (with $a,b$ stationary for $H$). Their $h$-variations contribute first at order four. In particular the coefficient $H$ at the seed is zero.

Let $k=x_2-\omega$. The reduced germ has no linear terms. Its only possible quadratic term is a multiple of $k^2$. Denote the homogeneous cubic part by
\[
P(k,h)=f k^3+g hk^2+e h^2k+jh^3.
\]
If the quadratic coefficient is nonzero, $\partial_kL=0$ forces $k=O(h^2)$ near the seed, and $\partial_hL=0$ then requires $j=0$. If the quadratic coefficient vanishes, a nonconstant stationary branch requires the two first derivatives of $P$ to have a common projective zero. Indeed write the Puiseux arc as $(k,h)=\tau^d(K,H)+o(\tau^d)$, where $d$ is the smaller of the two orders and $(K,H)\ne(0,0)$. In the gradient equations the terms of order $2d$ are $P_k(K,H)$ and $P_h(K,H)$. A nonzero discriminant excludes such a direction. For a binary cubic the common-projective-zero criterion follows by homogenizing the univariate derivative test, including the direction $h=0$, or by taking the homogeneous resultant of its partial derivatives. Thus it suffices to prove that $j$ and the discriminant are nonzero for generic $t$.

\subsection{The cubic at six equal masses}

Set $t=1$, so $m=2,n=3,S=1,M=6$, and write $u=p_{s5}^3=p_{s6}^3$, $r=p_{56}^3/u$, $q=(v_6 U_6)/(v_5 U_5)$. The binary equations and distance equations imply
\[
q^3=1,\quad 8(1+q)^3 r^2=1,\quad u=(1+2r)^{-1},
\quad U_5=\alpha,\quad U_6=-q\alpha,\quad \alpha^3=8/u^2.
\]
Indeed, at $t=1$ the centered binary has $v_5=-v_6=1/2$. Its large-coordinate equations are $v_i=S p_{si}^3 v_i+m_jp_{56}^3(v_i-v_j)$ for $\{i,j\}=\{5,6\}$. Dividing by $v_i$ gives $1=S p_{si}^3+mp_{56}^3$ in both rows. The center cubes therefore coincide, and at $t=1$ we have $1=u+2p_{56}^3$. The ratio of the squared lengths $D_{56}/D_{s5}=2(1+q)$ gives the equation for $r$. More generally, for positive $b_5=m_5,b_6=m_6$ the binary quantities satisfy $q^3=1$, $U_6=-q b_6 U_5/b_5$, $D_{56}/D_{s5}=m(b_5+q b_6)/(b_5 b_6)$, $r^2=(D_{56}/D_{s5})^{-3}$ and $1=(S+mr)u$. Since $D_{s5}=v_5 U_5$ and $p_{s5}^3=u$, the remaining scale satisfies $U_5^3=v_5^{-3}u^{-2}$. These equations give finite algebraic families. All denominators and internal distances stay nonzero above $t=1$. The $U$-Hessian at $1$, divided by the common nonzero factor $3p_{s5}/(4\alpha^2)$, is $\mathrm{diag}(1,q^{-1})+2r(1+q)^{-1}\left(\begin{smallmatrix}1&-1\\-1&1\end{smallmatrix}\right)$. Its determinant is $(1+2r)/q\ne0$.

At $t=1$ write $d=d_5$, $l=d_6/d_5$. Then $d^2=-2$ and $l^2=-1$.
It is convenient to compute the two outer Hessians in dimensionless
coordinates $Y_j=\beta y_j$. Put $p_0=(1+3\kappa)^{-1}$ and take both
values of $s_3$ in each row of the following table. Here $W_{ij}$ is
the signed potential on edge $ij$, divided by $\beta$, at the seed.
\[
\begin{array}{c|cccc}
 &s_3^2& \kappa& W_{js}&W_{ij}\ (i,j\ne s) \\ \hline
\text{equilateral} &3&1/(3s_3)&p_0&p_0/s_3\\
\text{collinear} &-3&1/(3s_3)&p_0/x_j&p_0\kappa(x_i-x_j)^2
\end{array}
\]
In both rows $\beta^3=(1+3\kappa)^2$.
Indeed $W_{js}=p_*r_j/\beta=p_0r_j$, and the collinear outer ratio
satisfies $q_{ij}=q_{ij}^3(x_i-x_j)^2$.
For an edge, put $A=\dot x_{ij}/x_{ij}$ and $B=\dot y_{ij}/y_{ij}$. The center has reference coordinates and variations zero. Its potential along a straight perturbation is its initial value times $(1+k A)^{-1/2}(1+k B)^{-1/2}$. This gives the edge contributions to the Hessian and to the third derivative explicitly.

For two variations $u,w$, the normalized Hessian is therefore
\begin{align*}
 Q_B(u,w)&=\frac12\sum_j(\dot x_j^u\dot y_j^w+\dot x_j^w\dot y_j^u)\\
 &\quad +\frac14\sum_{ij}W_{ij}\big[3(A_uA_w+B_uB_w)+A_uB_w+A_wB_u\big].
\end{align*}
The edge sum includes the three center edges. With $\dot x=(0,1,-1)$,
solve $Q_B((\dot x,\dot y),(0,e_j))=0$ for the three entries of $\dot y$.
This is precisely the derivative of transverse stationarity.
Set $b_1=\sum_j\dot y_j$ and
$c_1=\tfrac12\sum_j(W_{js}/x_j)(-3\dot x_j/x_j-\dot y_j/y_j)$.
These are the derivatives of $Y_B/\beta$ and $D_B/\beta$.
The coefficient $f$ is $\beta$ times the sum of the normalized edge values times $[-5(A^3+B^3)-3AB(A+B)]/16$. The quadratic inertia gives no straight-line cubic. Corrections from the curvature of the stationary manifold vanish by the differentiated small-coordinate stationarity equations and the choice of $x_2$ as parameter. Put $L_0=(1-q)r u$ and $A_0=(6L_0-1+q)/3$, so $\lambda=\alpha L_0$ and $a_0=\alpha A_0$. The remaining three cubic coefficients are
\[
g=-3d(1+l)/8,\qquad
e=\alpha\beta\big[(A_0-L_0)b_1-A_0c_1\big]/2,
\]
\[
j=(3d/2)\big[(1-lq-(1+l)A_0)\alpha+2(1-l)\beta\big].
\]
In particular $e$ is the derivative of $H$, with no contributions from varying the two stationary variables of $H$. The formula for $g$ uses $\sum_j\dot x_j^2/x_j=-1$.

The last expression for $j$ is the equilateral value. In the collinear
case the term $2(1-l)\beta$ is replaced by zero, because
$\sum x_j y_j=\sum x_j^2=0$. This follows directly by substituting
$g_{ij}=d_i x_j$ into $J$. The formulas for $g,e$ are the same for
both patterns when the indicated edge weights are used.

\subsection{A finite nonvanishing certificate}
The following finite-field enumeration includes all signs and both triangle orientations. For the equilateral case
use $p=73$, and for the collinear case use $p=97$.
\begin{enumerate}
\item In $\mathbb F_p$ loop over both primitive third roots $\omega$,
both indicated square roots $s_3$, all three values $q^3=1$, both
values $r^2=[8(1+q)^3]^{-1}$, both values $l^2=-1$ and both values
$d^2=-2$. All these polynomials split with the displayed root counts.
All denominators above are nonzero.
\item Compute $u=(1+2r)^{-1}$. Let $N$ be the least integer $a\ge2$
with $a^{(p-1)/3}\ne1$ in $\mathbb F_p$. In
$\mathbb F_p[\theta]/(\theta^3-N)$, loop over the three cube roots
$\alpha^3=8/u^2$ and the three cube roots $\beta^3=(1+3\kappa)^2$.
Each such root is $a\theta^i$ for a base-field $a$ and $0\le i<3$.
\item Form the binary determinant $(1+2r)/q$, the $3$-by-$3$ block
$Q_B((0,e_i),(0,e_j))$, and the four coefficients just given. Test the
two determinants, $j$, and
\[
 \Delta=g^2e^2-4fe^3-4g^3j-27f^2j^2+18fgej .
\]
\end{enumerate}
In each pattern the two transverse determinants, $j$ and $\Delta$ are nonzero on all $864$ tuples.
For example, in the collinear reduction $f=e=0$ and $\Delta=-4g^3j$.
The root counts and invertible denominators permit reduction of every
characteristic-zero tuple at a prime above $p$. A vanishing characteristic-zero
quantity would have given a vanishing residue. The binary equations, the
finite center and outer edge choices and the distant sign choices give finite
algebraic families regular above $t=1$. Locally their radical equations have
simple nonzero roots there. After discarding finitely many poles or vanishing
denominators, these families are finite covers of the mass line. Every
dominant component has a specialization above $1$, since the same radical
equations are monic with unit leading and constant coefficients near $1$.
The nonzero tests at all specializations show that the determinants, $j$
and $\Delta$ are nonzero away from finitely many mass values on every
dominant component. The reduced stationary germ has no singular branch on
that open set. This proves the required generic exclusion for both the
shared and binary placements.

\section{Diagram 94 with the unequal mass at an outer vertex}\label{nf:outer}
Relabel the quartet $0,1,2,3$ with $0$ marked, $3$ its shared center, and masses $(t,1,1,1)$. Take $X_0=1$, $X_3=0$. Summing \eqref{nf:forcex} gives $t+X_1+X_2=0$. Apply the torque equation \eqref{nf:torque}. For equal binary masses its coefficients $d_j$ satisfy $d_j^2=\pm 2i$. Put $x=v_1/v_0$ with $v_j=d_jm_j/\sqrt{X_j}$, so $v_2/v_0=-1-x$. Then
\[
 X_1=\frac{\delta_1}{t^2x^2},\qquad
 X_2=\frac{\delta_2}{t^2(1+x)^2},\qquad
 t^3x^2(1+x)^2+\delta_1(1+x)^2+\delta_2x^2=0,
\]
where $\delta_j=d_j^2/d_0^2\in\{1,-1\}$. After division by $t^3$ the quartic is monic over the local valuation ring and reduces to $x^2(1+x)^2$. Thus at infinity its roots tend to $0$ or $-1$. Exchange the two equal masses if necessary so that $x\to0$. Write $\nu(t)=-1$, $s=t^{-1/2}$, and $\delta=\delta_2$. Then
\[
 X_1\sim-t,\qquad X_2\sim\delta t^{-2},\qquad
 X_2=\frac{\delta s^4}{(1+x)^2},\qquad
 x=\frac{\lambda s^3}{\sqrt{1+s^2 X_2}},\qquad
 \lambda^2=-\delta_1.
\]
The square root in this last expression has constant term one. These relations determine a unique power series for each $(\delta,\lambda)$ by formal substitution.

\subsection{All possible valuations of the second projection}

Write $x_{ij}=X_i-X_j$, $y_{ij}=Y_i-Y_j$, $p_{ij}^2 x_{ij}y_{ij}=1$. All six differences in either projection and all six inverse distances are nonzero. Set $a_{ij}=\nu(x_{ij})$, $b_{ij}=\nu(y_{ij})$, $w_i=\nu(m_i)$. The six $a_{ij}$ are $(-1,0,0,-1,-1,2)$ in lexicographic edge order $01,02,03,12,13,23$.

Use the four additive triangle equations $y_{ij}-y_{ik}+y_{jk}=0$, the four $X$-force equations, and the three differences of the outer $Y$-force equations. An $X$-force from $j$ has doubled order $2w_j-a_{ij}-3b_{ij}$, and a $Y$-force from $j$ has doubled order $2w_j-3a_{ij}-b_{ij}$. In the difference of the $i,j$ outer $Y$ equations the single force on edge $ij$ has coefficient $m_i+m_j$. Each equation requires its minimum term order to occur at least twice. For example the doubled orders for the $X_0$ equation are
$0,1-3b_{01},-3b_{02},-3b_{03}$.
For the other rows the left sides contribute $2\nu(X_i)$ for $i=1,2$,
nothing for $X_3=0$, and $2b_{ij}$ for a $Y$ difference.
These eleven rational linear conditions force the following three vectors.

\[
\begin{array}{c|rrrrrr}
 &01&02&03&12&13&23\\ \hline
 A&0&-2&0&-2&1&-2\\ B&1/3&-2&-2&-2&-2&-2\\ C&0&0&-4&1&-4&-4
\end{array}
\]

The following exact rational linear feasibility test gives the table. For each
list of term orders $L$ append, for every indexed occurrence $u\in L$, the
clause $\bigvee_{v\in L\setminus\{u\}} v\le u$ (the removal deletes one
occurrence). Conjoin all clauses and the three conditions $b\ne b^{(A)}$,
$b\ne b^{(B)}$, $b\ne b^{(C)}$. The resulting rational linear query is
unsatisfiable under linear real quantifier elimination. For a direct implementation, the leading coefficients of
$(X_0,X_1,X_2)$ are $(1,-1,\delta)$ at respective orders $(0,-1,2)$.
The coefficient of an incoming negative $X$-force in a row brought to
the left side is $-\operatorname{lc}(x_{ij})P_{ij}^3$, and that of a
negative $Y$-force is $-P_{ij}/\operatorname{lc}(x_{ij})$. Here the
leading coefficient of each individual mass is $1$. For the direct
$ij$ interaction in a $Y$-difference multiply by
$\operatorname{lc}(m_i+m_j)$, equal to $2$ if $w_i=w_j$ and to $1$
otherwise. The other interactions in that difference retain their row
signs. These rules specify the coefficients as well as the orders. For each listed vector let $P_{ij}$ denote the nonzero leading coefficient of $p_{ij}$. The distance equation gives $\mathrm{lc}(y_{ij})=1/(\mathrm{lc}(x_{ij})P_{ij}^2)$. Retain the least-order terms in the same eleven equations and clear denominators. Take their Gröbner basis over $\mathbb Q$, adjoining the polynomial $1-H\prod P_{ij}$ to enforce nonzero edge coefficients. The reduced basis in graded reverse lexicographic order on $H,P_{01},\ldots,P_{23}$ is $[1]$ for both signs of $\delta$ in A and B and for $\delta=-1$ in C.

For C with $\delta=1$ the leading assignments in edge order are
\[
(P_{01},P_{02},P_{03},P_{12},P_{13},P_{23})=
(\epsilon A,A,-A,B,\eta A,A),
\qquad A^3=B^3=1,\quad \epsilon^2=1,\quad\eta^2=-1.
\]
Indeed, the $X_0,X_1$ forces give $P_{02}^3=P_{12}^3=1$ and the outer $Y$ difference $02$ gives $P_{23}=P_{02}$. Triangles $012,023,013$ give $P_{01}^2=P_{02}^2$, $P_{03}^2=P_{23}^2$, $P_{13}^2=-P_{03}^2$. The $X_3$ force fixes $P_{03}=-P_{23}$. These steps give exactly 36 initial assignments. Direct substitution makes every initial residual vanish.

\subsection{The shared-center force through order eight}

Use $s$ as uniformizing parameter in the coefficients. In case C write $p_{ij}=s^{k_{ij}}u_{ij}$ in lexicographic edge order, where
\[
(k_{ij})=(1,0,4,0,5,2).
\]
Keep the changes of variables exact by setting $D_{ij}=s^{-2a_{ij}}x_{ij}$ and $y_{ij}=s^{2b_{ij}}/(D_{ij}u_{ij}^2)$. Divide each force or additive relation by its least common power of $s$. Every resulting equation is regular at $s=0$ and the relevant denominators are units. Solve six equations consisting of triangles $012,013,023$, the forces $X_0,X_1$, and the $02$ outer $Y$-force difference. At each assignment their Jacobian in the six $u$'s is invertible. This follows successively from the powers, squares and linear equations in the previous paragraph, or by row reduction of the normalized six equations. Thus the formal implicit function theorem determines a single lift in ordinary powers of $s$, even if the original solution was sought in a further Puiseux extension.

For each of the 36 assignments and each of the four $\lambda\in\{1,-1,i,-i\}$, evaluate the normalized equations through order eight. If $E$ is the normalized six-component system, let $D=\partial_u E(0,u_0)$.
Starting from $u^{(<k)}=u_0+\cdots +s^{k-1}u_{k-1}$, take
$u_k=-D^{-1}[s^k]E(s,u^{(<k)})$ for $1\le k\le8$, then substitute the
result in all six equations. In the unused $X_3$ equation this gives the following first nonzero residuals.

\[
\begin{array}{c|rr}
 p&\text{first nonzero at }s^3&\text{first nonzero at }s^8\\ \hline
 97&108&36\\ 193&108&36
\end{array}
\]

These are first nonzero orders in the indicated residue fields. For instance, at 97 the $s^8$ coefficient is 50 in each of those 36 cases. The initial constants belong to $\mathbb Q(\zeta_{12})$, each prime splits there, and the enumeration includes all cube and fourth roots of unity. The Jacobians, leading differences and radical-expansion denominators used through order eight are invertible modulo each prime. Consequently the Taylor recurrence reduces coefficientwise. A nonzero residue proves that the corresponding characteristic-zero coefficient is nonzero. Either prime suffices.

Any solution over $\overline{\mathbb Q(t)}$ has a Puiseux expansion at a place over $t=\infty$ after extending constants. The preceding exhaustive argument rules out every such expansion, including those where original inverse distances or coordinates escape to infinity. Finally, the projection to the mass line of the algebraic necessary equations, with their nonvanishings imposed, is constructible. An infinite constructible subset of the line contains a nonempty open subset and has a nonempty generic fibre. The exceptional set is therefore finite. Together with the two lifting exclusions this proves Proposition~\ref{nf:prop}.
\section{The unit binary opposite a marked quartet}\label{ff:section}

For diagram 88 Chen and Tong obtain a necessary mass polynomial for a first limit on a curve through a positive physical point by taking a norm of the crossed torque over four-body shapes \cite[Prop.~5.14]{CT}. Suppose the mass $t$ belongs to the quartet in diagram 88. The two
masses of the separated binary are then one. We prove that the
necessary quartet equations have no point over
$\overline{\mathbb Q(t)}$. This step controls a complex algebraic
system, so it also applies to the physical subset of that system.
The choice of a finite test field is useful only after a boundary
analysis: the absence of solutions in a special affine fibre, by
itself, would allow the generic solutions to escape to infinity.

\paragraph{The equations to be excluded.}
Use isotropic coordinates $(z,w)$, with squared norm $zw$, and put
$m=(t,1,1,1)$, $M=\sum m_i$. For each of the six unordered pairs
introduce a signed inverse distance $p_{ij}=p_{ji}$ and set
$F^a_{ij}=p_{ij}^3(a_i-a_j)$ for $a=z,w$. In the normalized quartet
the force multiplier is one. The necessary equations of
\cite[Lemmas 6.16(a) and 6.17(a)]{Tomar} are
\begin{align}
 p_{ij}^2(z_i-z_j)(w_i-w_j)&=1,&
 \sum_{j\ne i}m_j F^a_{ij}&=a_i,\label{ff:quartet}\\
 z_i&\ne0,& \sum_i\frac{m_i\gamma_i}{\sqrt{z_i}}&=0,
 \qquad \gamma_i^4=1.\label{ff:torque}
\end{align}
Here $z$ denotes the large projection of the quartet. The coefficients
of the equal-binary torque are proportional to $1+\mathrm i$,
$1-\mathrm i$ and their negatives, giving exactly the fourth-root
phases after common rescaling. All square roots and all phases are
allowed. Summing the mass-weighted forces shows that both coordinates
are centered. The normalization $(z,w)\mapsto(cz,c^{-1}w)$ preserves
the system.

\subsection{Controlling specialization at the marked mass two}

Set $p=65521$ and $t=2+pU$. Give $\mathbb Q(U)$ the Gauss extension
of the $p$-adic valuation, so that the valuation of a polynomial is the
minimum of its coefficient valuations. The residue field is
$\mathbb F_p(\bar U)$, and every nonempty subsystem mass is a unit.
If a generic solution existed, extend the valuation to the finite
extension generated by its entries and square roots. By reciprocal
scaling of $z,w$, equalize the minimum valuation among positions and
forces in each color; denote it by $v_*$. A square root of a scaling
element suffices. Both minima cannot be positive. In fact, if the two
difference orders of one pair are $A,B$, its doubled force orders are
$-A-3B,-3A-B$. If $v_*=0$, all positions and forces are integral,
and these inequalities give $A=B=0$ on every pair.

Suppose $v_*<0$ and change the unit of valuation so $v_*=-2$. Circle
a vertex when its position has order $-2$ and mark an edge when its
force has order $-2$. The following list, up to labels and exchange of
the two colors, is valid also for this valuation. In the table $123$
in an edge column denotes the three edges of that triangle and $*$ in
an edge (respectively vertex) column denotes all six edges
(respectively four vertices).
\[
\begin{array}{c|cc|cc}
 &\multicolumn{2}{c|}{\text{circled vertices}}&
       \multicolumn{2}{c}{\text{marked edges}}\\
 &z&w&z&w\\ \hline
 \mathrm I&1,2&3,4&12&34\\
 \mathrm {II}&*&*&12,34&14,23\\
 \mathrm {III}&\varnothing&\varnothing&123&123\\
 \mathrm {IV}&*&\varnothing&123,14&123\\
 \mathrm V&\varnothing&\varnothing&*&*
\end{array}
\]
We give the finite linear check for this list because valuations at
a split prime need not obey sign restrictions on real coefficients.
Let $O$ be the origin. For each color use rational variables $A_{ij}$
or $B_{ij}$ for the orders on all pairs of the five symbols,
including $O$. An identically zero difference is omitted from every
sum: identify coincident labels in this construction. Since every
pair difference and every $z_i$ is nonzero, only two runs are needed,
with no identification and with $O$ identified with one $w$ vertex.
For a list $L$, let $E(L)$ say that its minimum is attained twice;
for the empty list put $E=\mathrm{true}$, for a singleton put
$E=\mathrm{false}$. This is the linear predicate
$\bigwedge_{u\in L}\bigvee_{v\in L,\,v\text{ at a different index}}v\le u$.
For each color, denote by $d_{ij}$ the coordinate order and by
$h_{ij}$ the force order ($2h_{ij}=-A_{ij}-3B_{ij}$ for $z$).
Use the following constraints.
\begin{enumerate}
\item For every triangle on the five symbols impose $E$ on its
three difference orders. Every $d_{iO}$ and $h_{ij}$ is at least
$-2$, and at least one attains $-2$ in each color.
\item For every nonempty subset $S$ of bodies and every anchor $a$
among the five symbols, impose $E$ on the concatenation of
$d_{aO}$, all $d_{ia}$ for $i\in S$, and all $h_{ij}$ for $i\in S$,
$j\notin S$. Omit the first entry for $a=O$. These are the terms in
the mass-weighted sum of the forces over $S$ about $a$, whose
coefficient of $a$ is the unit mass of $S$.
\item For every body pair $i,j$, impose $E$ on $d_{ij},h_{ij}$ and
all $h_{ik},h_{jk}$ for $k\notin\{i,j\}$. This is the difference
of the two force equations, with coefficient $m_i+m_j$ on the mutual
force.
\end{enumerate}
In each of the two runs conjoin the negations of the table patterns
for all $24$ labelings and both color orders. A pattern prescribes
equality to $-2$ at the specified entries and strict inequality at
all the others. Both resulting rational linear feasibility tests
return infeasible (quantifier-free linear real arithmetic). This
specifies all inputs of the support computation and uses only
necessary valuation conditions.

For the first four supports we give the initial-coefficient tests.
Write $\mu_i^2=m_i$ in the residue field, so the four masses, in
some order, are $2,1,1,1$. The four-body leading balances in
\cite[\S\S 5.1--5.4]{AK} can be used before their simplifications by
real positivity.
\begin{enumerate}
\item In I the two leading coordinates of each isolated pair are
opposite in the inverse ratio of its masses. Summing
$m_i q_i\wedge F_i$ on one pair eliminates its internal force. The
four interpair terms give the signed relation
\[
 \mu_1^3\mu_3^3\;\pm\;\mathrm i\mu_2^3\mu_3^3\;\pm\;
 \mathrm i\mu_1^3\mu_4^3\;\pm\;\mu_2^3\mu_4^3=0
\]
from equation (12) of that reference.
\item In II the two leading coordinates within each opposing cluster
agree. Eliminating their ratio from the two edge center equations
gives $m_1m_3=m_2m_4$, equation (14).
\item In III both colors balance internally on the triangle. In each
color the three cyclic force terms are proportional to the
opposite masses. Taking their quotients makes the two triples of
coordinate differences proportional. Closing the triangle gives
$1/\mu_1+e/\mu_2+f/\mu_3=0$ for $e^4=f^4=1$,
the relation immediately preceding equation (16).
\item In IV take the single edge in the $w$ color, with vertex 1
its attachment to triangle $123$. The three leading $w$ positions on
the triangle agree. Centering makes the fourth leading coefficient
different. Thus $w_{14},z_{14},p_{14}$ have orders $-2,2,0$, and
the three triangle differences have order one in each color. At
vertex 4 the $z$ forces have orders $2,5/2,5/2$, hence $z_4,z_1$
have order at least two. At triangle scale the three $z$ positions
are proportional to $(0,m_3,-m_2)$. With $R_{ij}$ the leading
signed distance, the internal forces and the closed $w$ triangle give
\[
 \frac{m_1}{R_{12}^3}=-\frac{m_2+m_3}{R_{23}^3}
       =\frac{m_1}{R_{13}^3},\qquad
 -\frac1{m_3 R_{12}}-\frac1{m_2 R_{13}}=\frac1{m_1 R_{23}} .
\]
Putting $\eta=R_{13}/R_{12}$ gives $\eta^3=1$ and
$m_1^2(m_2+m_3)(\eta m_2+m_3)^3=(m_2m_3)^3$.
\end{enumerate}
The field $\mathbb F_p$ splits the square roots of $2$ and $-1$
and all cube roots of unity. One obtains respectively $32,4,64,12$
nonzero differences on moving the right-hand sides to the left:
for each of the four possible marked positions take one square root
of each mass, all independent displayed signs in I, all $4^2$ phase
pairs in III and all three cube roots in IV. Thus all four supports
are impossible in the chosen residue characteristic.

In V both force orders attain $v_*$ on every edge, so
$A=B=-v_*/2$ and $v(p_{ij})=v_*/2$ everywhere. This leaves two
possibilities, bounded unit differences or this uniform collapse.
Divide $z,w$ by their respective differences on the same first pair,
and divide every inverse distance by $P$, the inverse distance of
that pair. The distance identities persist and the common force
multiplier becomes $P^{-3}$. The new differences and signed inverse
distances are units. The centered positions are integral, since
$a_i=M^{-1}\sum_jm_j(a_i-a_j)$. At most one high coordinate can
have positive valuation. If one did, its torque term would be the
only term of negative valuation, contrary to \eqref{ff:torque}.
Thus all four high coordinates are units. Reduction gives a finite
projective quartet in characteristic $p$, with masses $(2,1,1,1)$,
all high coordinates nonzero, and multiplier $\lambda$, possibly
zero. All the following tests have coefficients in $\mathbb F_p$.
Emptiness over its algebraic closure implies emptiness over every
extension, including extensions of $\mathbb F_p(\bar U)$, by the
Nullstellensatz after adjoining inverses for the nonvanishing conditions.

\subsection{The torque polynomial and aligned strata}

For centered high coordinates $H_i$ put
$A_i=e_i m_i^2\prod_{j\ne i}H_j$, where $e_0=1$ and the three
other $e_i$ range independently over $\{1,-1\}$. Define
\begin{equation}
 \Phi(A)=\big[(A_0+A_1-A_2-A_3)^2-4A_0A_1-4A_2A_3\big]^2
              -64A_0A_1A_2A_3 .\label{ff:phi}
\end{equation}
If the four torque terms $d_i$ add to zero, first square
$d_0+d_1=-(d_2+d_3)$ and then eliminate the two remaining products
by squaring. The result is $\Phi(d_0^2,d_1^2,d_2^2,d_3^2)=0$.
The $A_i$ are proportional to these squares for one of the eight
patterns. Only necessity of this polynomial test is used.

If all four points lie on a line, its direction is nonisotropic
because the six squared distances are nonzero. Normalize the three
unit positions and the marked position to $q=(0,1,x,y)$ in that
order. For every pair $i<j$ choose $\sigma_{ij}=\pm1$ and add
$m_j\sigma_{ij}/(q_i-q_j)^2$ to $f_i$, subtracting
$m_i\sigma_{ij}/(q_i-q_j)^2$ from $f_j$. These are the signed
accelerations, with a common scale absorbed. They must be affine in
$q_i$ even if $\lambda=0$. Put
\[
 D=\prod_{i<j}(q_i-q_j)^2,\qquad
 G_i=D\{f_i-f_0-q_i(f_1-f_0)\},\quad i=2,3 .
\]
The following univariate elimination takes place over $\mathbb F_p$.
For each sign choice form $\operatorname{Res}_y(G_2,G_3)$, always
nonzero. For each distinct irreducible factor other than $x,x-1$,
specialize $x$ to its residue class in the field defined by that
factor and compute $g=\gcd_y(G_2,G_3)$. Remove all factors of
$y(y-1)(y-x)\prod H_i$, repeatedly taking gcd and dividing, where
$H_i=5q_i-(1+x+2y)$. For each of the eight signs
$\gcd(g,\Phi(A))=1$. As intermediate checks, enumerate signs in
lexicographic order with $+1$ first, pairs in increasing label order;
skip previously seen orbits of the three unit-body permutations and
simultaneous sign reversal, reversing a sign when the permuted edge
order reverses. The six sums of $\deg_x(\mathrm{factor})\deg_y g$
before applying the torque are $35,35,35,35,33,35$. Thus this test
also covers zero-multiplier collinear configurations.

If exactly three positions are collinear, choose two unit masses on
their line at $0,1$, place the third at $x$ with mass $m=1$ or $2$,
and the off-line point at $(v,h)$ with $h\ne0$. These are orthogonal
coordinates over the algebraic closure. The transverse forces at
the three aligned points show that their three inverse cubes to the
off-line point equal $k_0=\lambda/5$. Thus $\lambda\ne0$ and their
squared distances are $S,aS,zS$ with $a^3=z^3=1$. Subtracting their
squared lengths gives $1-2v=(a-1)S$ and
$S=x(x-1)/[(z-1)-x(a-1)]$. The denominator is nonzero. Fixing the
first internal edge sign, the internal inward accelerations are
\[
 f_0=1+mb/x^2,\qquad f_1=-1+mc/(1-x)^2,\qquad
 f_2=-b/x^2-c/(1-x)^2,\qquad b,c=\pm1 .
\]
Subtracting axial forces gives $f_2-f_0=x(f_1-f_0)$ and
$f_1-f_0=-(2+m)k_0$. Hence the two rational functions
$f_2-f_0-x(f_1-f_0)$ and $(f_1-f_0)^2 S^3-(2+m)^2$ must vanish.
For all $2\cdot 2^2(3^2-1)=64$ choices other than $(a,z)=(1,1)$,
take their numerator gcd. Removing all factors of
$x(x-1)[(z-1)-x(a-1)]\operatorname{num}(f_1-f_0)$ leaves a
constant in every case. This excludes the stratum by forces alone.

\subsection{Stress reduction for genuinely planar quartets}

We may now assume that no three points are collinear. Put
$k=\lambda/5$. Form the symmetric matrix $S$ with off-diagonal
entries $m_i m_j(p_{ij}^3-k)$ and diagonal entries chosen for zero
row sums. Its kernel contains the independent vectors
$\mathbf1,z,w$ by the force equations. If $S=0$ then $k\ne0$ and
all inverse cubes are equal. Normalize the first squared distance
to one. Every squared distance then belongs to $\mu_3$, the three
cube roots of unity. Write $D_{ij}$ for these squared distances,
$a=D_{02}$, $b=D_{03}$, $U=1+a-D_{12}$, $V=1+b-D_{13}$,
$W=a+b-D_{23}$. A doubled Gram determinant on the three vectors
from vertex zero is
\[
 8ab+2UVW-2W^2-2aV^2-2bU^2 .
\]
It must vanish in a plane. For the $3^5$ tuples in $\mu_3^5$ it
takes 23 distinct residues modulo $65521$, all nonzero.

If $S\ne0$, write $S=c vv^{\mathsf T}$, $c\ne0$. The entries of
$v$ give the unique affine dependence and none vanishes, since no
three points are collinear. Translate the marked point to the origin,
call the other points $Q_1,Q_2,Q_3$, and put $a_i=p_{0i}^3-k$.
Indeed, writing $T_i=v_i/m_i$, we have $a_i=cT_0T_i$, so with
$L=(cT_0^2)^{-1}$ the off-diagonal entries and the affine dependence
give
\begin{equation}
  p_{ij}^3=k+L a_i a_j,\qquad
 2=-L\sum_{i=1}^3a_i,\qquad \sum_{i=1}^3a_i Q_i=0 .
 \label{ff:stress}
\end{equation}
Here $L$ and every $a_i$ are nonzero. The marked mass enters the
second equation through $v_0+v_1+v_2+v_3=0$.

\paragraph{The zero multiplier.}
Aquino-Michaels confines masses supporting a zero-force quartet to a proper algebraic subset by treating signed collinear forces and noncollinear rank-one stresses \cite[Lemmas~5.3--5.5, Thm.~5.6]{Residual}. We check the relevant special fibre here.
When $k=0$, multiply all inverse distances by $L^{1/3}$ and divide
the positions by that scalar. Write $b_i$ for the new star inverse
distances. Equation \eqref{ff:stress} becomes
$p_{ij}^3=b_i^3b_j^3$, $\sum b_i^3=-2$ and
$\sum b_i^3Q_i=0$. Use the bilinear Euclidean square and set
$V_i=b_i^2Q_i$. Since $Q_i^2=b_i^{-2}$, direct expansion gives
\[
 V_i^2=b_i^2,\qquad
 \eta_{ij}=(V_i-V_j)^2=b_i^2 b_j^2(Q_i-Q_j)^2,\qquad
 \eta_{ij}^3=1\quad(i\ne j).
\]
Put $\eta_{jj}=0$. Dot $\sum b_i V_i=0$ with $2V_j$, and use
$2V_i\cdot V_j=V_i^2+V_j^2-\eta_{ij}$ to obtain four cubics
\[
 (b_1+b_2+b_3)b_j^2-2-\sum_i\eta_{ij}b_i=0\quad(1\le j\le3),
 \qquad b_1^3+b_2^3+b_3^3+2=0 .
\]
For each of the ten multisets of size three from $\mu_3$, assign
the three roots to edges $12,13,23$. Every ordering is equivalent
under permutations of the three outer points. The reduced
degree-reverse-lexicographic Gr\"obner basis of these four cubics over
$\mathbb F_{65521}$ is $\{1\}$ in each case. These are all inputs
needed for the zero-multiplier test. This step explains why it was
legitimate to allow a uniform collapse in the valuation analysis.

\subsection{Three finite quotient algebras for the nonzero multiplier}
\label{ff:charts}

For $k\ne0$ multiply inverse distances by $k^{-1/3}$ and positions
by its reciprocal. The offset in \eqref{ff:stress} is now one. Let
$b_i$ be the normalized inverse star distances, $\alpha_i=b_i^3-1$
and $T=\sum\alpha_i$. Equation \eqref{ff:stress} gives
$T\ne0$, $\sum\alpha_iQ_i=0$ and
$p_{ij}^3=1-2\alpha_i\alpha_j/T$. Choose a nonzero scale $B$,
write $u_i=b_i/B$, $s=B^3$, and use positions $BQ_i$. Their squared
star lengths are $R_i=u_i^{-2}$. Squaring
$\alpha_i BQ_i+\alpha_j BQ_j=-\alpha_k BQ_k$ yields the squared
outer distance
\begin{equation}
 D_{ij}=\frac{(\alpha_i+\alpha_j)(\alpha_i R_i+\alpha_j R_j)
             -\alpha_k^2R_k}{\alpha_i\alpha_j}.
 \label{ff:distance}
\end{equation}
The signed outer inverse cube then implies
$(T-2\alpha_i\alpha_j)^2D_{ij}^3=T^2 s^2$.
For $g=u_1u_2u_3$ put
\begin{align*}
 N_{ij}&=(\alpha_i+\alpha_j)(\alpha_i u_j^2+\alpha_j u_i^2)
            u_k^2-\alpha_k^2u_i^2u_j^2,\\
 G_{ij}&=(T-2\alpha_i\alpha_j)^2N_{ij}^3
                   -T^2s^2(g^2\alpha_i\alpha_j)^3 .
\end{align*}
Here $\alpha_i=s u_i^3-1$. We give explicit generators for the
three finite rings; all polynomial calculations henceforth are modulo
$65521$.

\paragraph{Distinct ratios with nonzero sum.}
Put $B=\sum b_i$, so $\sum u_i=1$. Write $f=\sum_{i<j}u_i u_j$
and $Q(Z)=Z^3-Z^2+fZ-g$. At $Z=u_k$ define
\[
 U=1-Z,\quad V=Z^2-Z+f,\quad C=sZ^3-1,\quad
 T=s(1-3f+3g)-3 .
\]
Since $U=u_i+u_j$ and $V=u_i u_j$, the opposite edge has
\[
 J=\alpha_i\alpha_j=s^2V^3-s(U^3-3UV)+1,\qquad
 N=(T-C)\{sUg^2-(U^2-2V)Z^2\}-C^2V^2 .
\]
In $\mathbb F_p[v,s,f,g]$ take as generators the three coefficients
of the remainder of $(T-2J)^2N^3-T^2s^2g^6J^3$ upon division by
the monic $Q$, and the fourth generator
$vsgT\operatorname{disc}(Q)\operatorname{Res}_Z(Q,C)-1$.
Vanishing of the remainder follows from the three distinct roots;
the inverse variable $v$ implements precisely the indicated open
conditions.

\paragraph{Zero sum and repeated ratios.}
In both rings divide by $B=b_1$. For zero sum take
$(u_1,u_2,u_3)=(1,x,-1-x)$. For a repeated ratio relabel the
equal pair and take $(1,1,x)$. Use in $\mathbb F_p[v,x,s]$ the
three generators $G_{12},G_{13},G_{23}$ and
$vsgT\alpha_1\alpha_2\alpha_3-1$. The three charts cover all
ratios including repeated cubes with distinct inverse distances.

\paragraph{Computable finite-algebra checks.}
Here are the intermediate sizes. In the displayed variable orders
use degree reverse lexicographic order. The middle column counts
minimal generators of the leading monomial ideal; the next counts
monomials outside it.
\[
\begin{array}{l|r|r|l}
 \text{chart}&\#\text{ leading generators}&
 \text{quotient dimension }n&\text{separating element }\ell\\ \hline
 \text{distinct, nonzero sum}&165&420&g\\
 \text{zero sum}&3&4&-(v+x+s)\\
 \text{repeated}&36&102&s
\end{array}
\]
For clarity we describe how to reproduce the complete test without
receiving a list of candidate configurations. For each quotient
compute a Gr\"obner basis of the displayed generators. Start at the
constant monomial $1$ and repeatedly multiply by each variable,
retaining a new monomial if none of the leading monomials divides it.
This gives a basis $\mathcal B$ of the stated size. Define $M_h$
by reducing $h b$ for every $b\in\mathcal B$ and placing the
resulting coefficient lists in columns. At this point only exact
linear algebra is needed:
\begin{enumerate}
\item Compute $P(Y)=\det(YI-M_\ell)$. The gcd of $P$ with $P'$
is one in every chart. In the largest chart the irreducible factor
degrees are $2,5,17,156,240$. Because $P$ is squarefree of degree
$n$, the powers $1,\ell,\ldots,\ell^{n-1}$ are a basis.
Writing $e$ for the column of the constant monomial, the coordinate
polynomial $H_h(Y)$ of each variable $h$ is obtained by solving
\[
 [e,M_\ell e,\ldots,M_\ell^{n-1}e]\,
     \operatorname{coeff}(H_h)=M_h e .
\]
Equivalently one can verify a squarefree degree-$n$ parametrization
by substituting its coordinates into every defining generator and
checking that $\ell$ equals the parameter. Those checks certify $n$
distinct geometric points in a quotient of dimension $n$.
\item For every irreducible factor $P_0$ of $P$, put
$K=\mathbb F_p[Y]/(P_0)$ and evaluate $H_h$ there. In the first
chart adjoin $u_1$ by the monic cubic $Q$, then $u_2$ by
$u_2^2-(1-u_1)u_2+u_1^2-u_1+f=0$, and put
$u_3=1-u_1-u_2$. The rank-six algebra includes every ordering.
In the two ordered charts use their displayed ratios directly.
\item Normalize the relative high coordinates to $(0,1,h,H_3)$.
The low coordinates of points 1 and 2 relative to point zero are
then $R_1,R_2/h$, so their squared separation gives the monic
equation
\[
 h^2-\frac{R_1+R_2-D_{12}}{R_1}h+\frac{R_2}{R_1}=0,\qquad
 H_3=-\frac{\alpha_1+\alpha_2 h}{\alpha_3}.
\]
Use \eqref{ff:distance} for $D_{12}$ and adjoin $h$ by this
quadratic. For $Z=(0,1,h,H_3)$, put $H_i=5Z_i-(1+h+H_3)$.
These are the centered coordinates multiplied by five.
For every one of the eight signs compute \eqref{ff:phi} with
$A_i=e_i m_i^2\prod_{j\ne i}H_j$.
\item Descend through the quadratic, and through the ordering
algebras when present, replacing each value by the determinant of
multiplication by it. Explicitly in a monic extension
$R[y]/(C)$ of degree $d$ use as columns the coefficients of the
remainders of $a,ya,\ldots,y^{d-1}a$ modulo $C$. The eight
resulting elements of $K$ are nonzero for every $P_0$ in all three
charts. A nonzero final determinant makes $\Phi$ a unit, even if a
quadratic above a particular base point has a double root.
\end{enumerate}
Thus neither projection, for any ordering or signs, can realize the
torque at a regular residue point. Together with the aligned,
zero-stress, and zero-multiplier cases, this exhausts the proposed
finite fibre. The specialization argument now excludes
\eqref{ff:quartet}--\eqref{ff:torque} over
$\overline{\mathbb Q(t)}$. Equivalently the constructible image
of this necessary signed system on the mass line has finite support:
an infinite constructible subset of a line contains a nonempty open
set and would have a geometric generic point.
\section{Conclusion and verification}\label{end:proof}
\begin{proof}[Proof of Theorem~\ref{intro:main}]
Fix a positive real ratio $t$ outside the union of the exceptional values in the preceding reductions. Suppose that there are infinitely many real shapes. By Sections~\ref{intro:places} and~\ref{br-real}, at this fixed $t$ there is a curve on the normalized complex slice with a nonconstant squared distance and a singular place satisfying $v(\Delta)\le0$. The two origin leaves agree. The quantities $E$ and $H$ have the constant values of a real configuration along this curve.

Sections~\ref{early:census} and~\ref{early:98} specialize the catalogue, exclude the two-triangle diagram and produce 852 marked placements. Section~\ref{br-ledger} leaves 187 placements in three diagrams. Among these, 90 placements in 29 symmetry classes satisfy the origin condition. Section~\ref{br-real} uses the energy and its derivative to exclude diagram 93 and the marked binary in diagram 88. The possible place now belongs to one of 74 placements in 22 classes. In 54 placements the diagram is 94. Proposition~\ref{nf:prop} and its proof in Sections~\ref{nf:signed}--\ref{nf:outer} exclude all three roles of the marked mass. In the other 20 placements the diagram is 88 and the quartet contains the marked mass. Section~\ref{ff:section} excludes its force and torque system. There is therefore no place satisfying the necessary conditions.

To check the parameter quantifier, the original catalogue has finitely many cells. The nonzero specialized mass obstructions give finitely many roots. The branch degeneracies, denominator zeros and losses of constant degree in the finite algebraic tests also have finite algebraic support. Where the exclusion is first made over the geometric generic field, it is applied to a finite polynomial system with its listed inequations, as in Section~\ref{intro:generic}. For the physical inequalities, the extra omissions consist of the vertical images, poles and branch values of the finite algebraic incidences described in Section~\ref{br-real}. All these constructions are over the algebraic numbers. Intersecting their finite union with the positive real axis gives a permissible $S$.
\end{proof}

\paragraph{What an exceptional ratio would mean.}
The proof exhibits mechanisms by which the generic exclusions might fail at a special ratio. The vanishing of an obstruction is only a necessary condition for the corresponding end. Even the existence of a signed singular end would leave the separate question of a real family at that ratio. Thus the argument gives neither an example with infinitely many physical shapes nor an enumeration of shapes. At $t=1$ the existing equal-mass enumeration already supplies finiteness \cite{MZplanar}.

\paragraph{Reproducing the finite steps.}
The separate \suppfile{} contains the verifier sources, the input conventions for the external catalogue, and instructions for replaying the following calculations. The analytic arguments, normalized equations and specialization steps are given in the text.
\begin{itemize}
\item To reconstruct the catalogue, run \texttt{python3 run.py catalogue} with \texttt{--source}, \texttt{--cores}, and \texttt{--hand} pointing to the source tree and the two certificate archives identified in the supplement \cite{Data}. It specializes the core mass identities and compares the resulting lists with the recorded tables. The lists distinguish a marked placement (one cell and one choice of mass slot) from its symmetry class.
\item Run \texttt{python3 run.py arithmetic} for the diagram tests. The triangle tests use exact rational enclosures. Other tests include polynomial greatest common divisors, finite-field calculations, exhaustive constraint checks, and the polynomial-system solver \texttt{msolve}~\cite{msolve} for the three regular quartet charts. The detailed equations in Sections~\ref{nf:signed}--\ref{ff:section} specify which denominators may be inverted and which other charts cover their zeros.
\item The supplement also identifies the Lean statements for the algebraic identities, inequalities or finite enumerations that have been formalized. The physical-to-algebraic reduction, catalogue application, lifting arguments and the entire implication in Theorem~\ref{intro:main} have a mathematical exposition here. The end-to-end theorem has no Lean formalization in this release.
\end{itemize}

\paragraph{Acknowledgment and provenance.}
The catalogue and the diagrams are due to the external sources listed below. This article incorporates and expands arguments developed in the author's earlier manuscripts on the five-equal-mass family, the reduction of its diagram list, the physical energy bounds and the final two diagrams. All four stages used in the proof are presented here. The supplementary programs give machine-readable forms of its finite computations.

\begin{thebibliography}{99}\small\setlength{\itemsep}{1pt}
\bibitem{Smale} S. Smale. Mathematical problems for the next century. \emph{The Mathematical Intelligencer} \textbf{20} (1998), 7--15. \href{https://doi.org/10.1007/BF03025291}{doi:10.1007/BF03025291}.
\bibitem{AK} A. Albouy and V. Kaloshin. Finiteness of central configurations of five bodies in the plane. \emph{Annals of Mathematics} \textbf{176} (2012), 535--588. \href{https://doi.org/10.4007/annals.2012.176.1.10}{doi:10.4007/annals.2012.176.1.10}.
\bibitem{CC} K.-M. Chang and K.-C. Chen. Toward finiteness of central configurations for the planar six-body problem by symbolic computations. \href{https://arxiv.org/abs/2303.02853}{arXiv:2303.02853} (2023). The journal articles are \href{https://doi.org/10.1016/j.jsc.2023.102277}{(I) Determine diagrams and orders}, \emph{J. Symbolic Comput.} \textbf{123} (2024), 102277, and \href{https://doi.org/10.1137/24M1716070}{(II) Determine mass relations}, \emph{SIAM J. Appl. Dyn. Syst.} \textbf{24} (2025), 2369--2404.
\bibitem{CT} Z. Chen and Z. Tong. Finiteness of planar central configurations of six bodies for all positive masses: Part I, exceptional-mass first limits. \href{https://arxiv.org/abs/2610.23329v1}{arXiv:2610.23329v1} (2026).
\bibitem{Residual} K. Aquino-Michaels. Reducing the Chang--Chen residual: certified eliminations in the planar six-body central-configuration problem. \href{https://github.com/no-way-labs/smale6-residual/tree/07f6c99}{Preprint and certificates, revision 07f6c99} (2026).
\bibitem{GenericN} K. Aquino-Michaels. Generic finiteness of planar central configurations for every number of bodies. \href{https://doi.org/10.5281/zenodo.24371245}{Preprint, Zenodo 24371245} (2026).
\bibitem{DP} T. Dias and B.-Y. Pan. Generic finiteness for a class of symmetric planar central configurations of the six-body problem and the six-vortex problem. \emph{J. Dynam. Differential Equations} \textbf{32} (2020), 1579--1602. \href{https://doi.org/10.1007/s10884-019-09765-y}{doi:10.1007/s10884-019-09765-y}.
\bibitem{msolve} J. Berthomieu, C. Eder and M. Safey El Din. msolve: a library for solving polynomial systems. \emph{Proceedings of ISSAC 2021}, 51--58. \href{https://doi.org/10.1145/3452143.3465545}{doi:10.1145/3452143.3465545}.
\bibitem{MZplanar} M. Moczurad and P. Zgliczy\'nski. Central configurations in planar $n$-body problem with equal masses for $n=5,6,7$. \emph{Celestial Mechanics and Dynamical Astronomy} \textbf{131} (2019), article 46. \href{https://doi.org/10.1007/s10569-019-9920-6}{doi:10.1007/s10569-019-9920-6}.
\bibitem{MZnear} M. Moczurad and P. Zgliczy\'nski. Central configurations in the spatial $n$-body problem for $n=5,6$ with equal masses. \emph{Celestial Mechanics and Dynamical Astronomy} \textbf{132} (2020), article 56. \href{https://doi.org/10.1007/s10569-020-09993-1}{doi:10.1007/s10569-020-09993-1}.
\bibitem{Tomar} T. S. Tomar. Smale's sixth problem for six bodies in the plane: finiteness of central configurations for generic masses, with a proof formalized in Lean. Research preprint, \href{https://hexagonmath.org/pdf/2611.00377v1}{Hexagon 2611.00377v1} (2026).
\bibitem{Data} T. S. Tomar. Computations and certificates for the six-body finiteness theorem. \href{https://github.com/tejasvi/smale6-six-bodies-generic/tree/84bbcf276d029fdf31375c4a21c6a7e136955e51}{Catalogue and source, revision 84bbcf2}; \href{https://zenodo.org/records/23472605}{Zenodo record 23472605} (2026).
\end{thebibliography}
\noindent\small\copyright\ 2026 Chatgpt. This paper and its \LaTeX{} source are licensed under \href{https://creativecommons.org/licenses/by/4.0/}{Creative Commons Attribution 4.0 International}.
\end{document}
